\documentclass[aps,preprint]{revtex4}%
\usepackage{amsfonts}
\usepackage{amsmath}
\usepackage{amssymb}
\usepackage{graphicx}%
\setcounter{MaxMatrixCols}{30}
%TCIDATA{OutputFilter=latex2.dll}
%TCIDATA{Version=4.00.0.2321}
%TCIDATA{CSTFile=revtex4.cst}
%TCIDATA{Created=Monday, September 20, 2004 09:58:21}
%TCIDATA{LastRevised=Saturday, December 31, 2005 11:37:47}
%TCIDATA{<META NAME="GraphicsSave" CONTENT="32">}
%TCIDATA{<META NAME="DocumentShell" CONTENT="Articles\SW\REVTeX 4">}
%TCIDATA{Language=American English}
\newtheorem{theorem}{Theorem}
\newtheorem{acknowledgement}[theorem]{Acknowledgement}
\newtheorem{algorithm}[theorem]{Algorithm}
\newtheorem{axiom}[theorem]{Axiom}
\newtheorem{claim}[theorem]{Claim}
\newtheorem{conclusion}[theorem]{Conclusion}
\newtheorem{condition}[theorem]{Condition}
\newtheorem{conjecture}[theorem]{Conjecture}
\newtheorem{corollary}[theorem]{Corollary}
\newtheorem{criterion}[theorem]{Criterion}
\newtheorem{definition}[theorem]{Definition}
\newtheorem{example}[theorem]{Example}
\newtheorem{exercise}[theorem]{Exercise}
\newtheorem{lemma}[theorem]{Lemma}
\newtheorem{notation}[theorem]{Notation}
\newtheorem{problem}[theorem]{Problem}
\newtheorem{proposition}[theorem]{Proposition}
\newtheorem{remark}[theorem]{Remark}
\newtheorem{solution}[theorem]{Solution}
\newtheorem{summary}[theorem]{Summary}
\newenvironment{proof}[1][Proof]{\noindent\textbf{#1.} }{\ \rule{0.5em}{0.5em}}
\begin{document}
\preprint{ }
\title[Janaks Theorem]{Janaks Theorem and Energy Functionals of Orbitals and Occupation Numbers}
\author{B. Weiner}
\affiliation{Pennsylvania State University}
\keywords{}
\pacs{}

\begin{abstract}
Janaks theorem shows that the rate of change of the total N-electron energy
with respect to the occupation numbers of the natural orbitals of the model
Kohn-Sham (KS) independent particle state are equal to the Lagrange
multipliers associated with the normalization constraint imposed on these
orbitals. However many conceptual questions surround this theorem. In what
sense can the occupation numbers vary as the model KS state is by definition
an Independent Particle State (IPS) with integer occupation numbers? What is
the functional relationship between the eigenvalues and the occupation number?
How does one define a KS functional that is differentiable with respect to
occupation numbers? In order to examine the meaning and content of this
theorem we review the construction of the implicitly defined Hohenberg-Kohn
(HK) and KS energy functionals and obtain an expression of Janaks theorem in
the context of constrained optimization theory. It is shown that, in fact,
Janaks theorem holds for any total energy functional defined in terms of
orbitals and occupation numbers. In particular it applies to Antisymmetrized
Geminal Power energy functionals, which also solely depend on orbitals and
occupation numbers and whose functional form is explicitly known. The
properties and structure of these orbital functionals are contrasted and
compared and subtle differences are discussed..\qquad

\end{abstract}
\date{09/07/2005}
\startpage{1}
\maketitle
\tableofcontents


\section{ Introduction\label{Intro}}

Janaks theorem is a useful auxiliary tool in the Kuhn-Sham (KS) approach to
Density Functional Theory (DFT), that is helpful in analyzing the structure
and meaning of the eigenvalues of the KS Fock operator. It states that these
eigenvalues are equal to the derivatives of the exact total energy with
respect to the occupation numbers of the First Order Reduced Density Operator
(FORDO) of the model KS state. However many conceptual questions surround this
theorem. In what sense can the occupation numbers vary as the model KS state
is by definition an Independent Particle State (IPS)\ with integer occupation
numbers? What is the functional relationship between the eigenvalues and the
occupation number? How does one define a KS\ functional that is differentiable
with respect to occupation numbers?

In both KS-DFT and Hartree-Fock theory these eigenvalues correspond to
Lagrange parameters associated with the normalization constraint on the
orbitals, that constitute the Independent Particle State (IPS) and to
ionization energies of electrons described by these orbitals. This
correspondence can be generalized to Multi Configurational Self Consistent
States (MCSCF), as described for instance in \cite{WeinerOrtiz05}. Constrained
optimization theory \cite{Hestenes} shows, if the Hessian of the Lagrangian is
nonsingular, that the values of Lagrange parameters correspond to the
sensitivity of the object function to the control parameters, which are the
values of the constraints. In this case the $N$-electron energy to variations
of the normalization of the orbitals. Hence Janaks theorem implies that the
normalization of the orbitals must be a function of the occupation numbers.
The form of the dependence of the energy on the occupation numbers can be
factored through the dependency of the energy on the orbitals, which in
particular contains the dependence on the norm of the orbitals, and thence on
the occupation numbers. This factoriztion is especially important as the
relationship between the form of the orbitals and the total energy is crucial
in the construction of molecular models and hence in the design of new
materials and molecules.

One can construct many different types of orbital functionals, we will
consider the following ones

\begin{itemize}
\item Functionals that are invariant to local phase changes of orbitals thus
preserve the electron density

\item Functionals that are invariant to overall phase changes of the orbitals
i.e. to the group $U\left(  1\right)  $, e.g. Density Matrix and KS\ Functionals

\item Functionals that are not invariant to $U\left(  1\right)  $ e.g.
correlated AGP\ functionals.
\end{itemize}

In order to answer the preceding questions and analyze the energy dependence
on occupation numbers we will review the basic principles of Hohenburg-Kohn
(HK), KS DFT (section 2), Janaks Theorem (section 3) and Density Matrix
Functional Theory (DMFT) in the context of a constrained variational approach
and sensitivity theory. 

The KS\ density functional, density matrix functional and the AGP functional
determine the total energy in terms of a set of occupations numbers and one
particle orbitals. The first two have only been defined implicitly while the
AGP functional has been explicitly constructed \cite{Weiner&Ortiz2005} in
terms of General Spin Orbitals (GSO's) and parameters that determine the FORDO
occupation numbers (see section 4). In section 5 we investigate the form of
the derivatives of this functional with respect to occupation numbers and the
special case of its value at integer occupation. In section 6 we conclude with
a discussion and summary.

\section{Density Functionals \label{Standard}}

In this section we review some fundamental background needed to defined HK and
KS density functionals. We will mention important technical mathematical
aspects that need to be examined, but defer a full analysis of these issues to
a future article and we will focus for simplicity on the mathematically most
well behaved cases.

\subsection{Hohenberg-Kohn DFT}

The HK density functional is defined in terms of the kinetic energy operator,
$K,$ the coulomb operator, $V_{2},$ and the charge density $\rho$.

\subsubsection{Densities and States}

\bigskip Density Functional Theory is based on the existence of a universal
functional, $\mathcal{\tilde{F}}\left(  \rho\right)  $ \cite{HK} of the charge
density $\rho\left(  \mathbf{r}\right)  ,$ where
\begin{equation}
\rho\left(  \mathbf{r}\right)  =\int D^{1}\left(  \mathbf{r,\xi;r,\xi}\right)
d\mathbf{\xi}%
\end{equation}
and $D^{1}\left(  \mathbf{r,\xi;r}^{\prime}\mathbf{,\xi}^{\prime}\right)  $ is
the integral kernel of the First Order Reduced Operator (FORDO) $D^{1}$
\cite{Coleman} (see Appendix \ref{DenMat} for details of density and reduced
density operators and matrices), which can be expanded as
\begin{equation}
D^{1}=\int D^{1}\left(  \mathbf{r,\xi;r}^{\prime}\mathbf{,\xi}^{\prime
}\right)  \left\vert \mathbf{r,\xi}\right\rangle \left\langle \mathbf{r}%
^{\prime}\mathbf{,\xi}^{\prime}\right\vert d\mathbf{r}d\mathbf{r}^{\prime
}d\mathbf{\xi}d\mathbf{\xi}^{\prime}%
\end{equation}
in the basis $\left\{  \left\vert \mathbf{r,\xi}\right\rangle \right\}  $
formed by the tensor product of the generalized eigenvectors, $\left\{
\left\vert \mathbf{r}\right\rangle \right\}  ,$ of the position operator and a
family of $SU(2)$-coherent spin states $\left\{  \left\vert \mathbf{\xi
}\right\rangle \right\}  $ $\cite{weiner2005}$. In order to consider spin
dependent effects it is convenient to extend the original formulation of DFT
in terms of $\rho\left(  \mathbf{r}\right)  $ and consider a universal
functional, $\mathcal{F}\left(  \gamma\right)  $, of the space-spin density
$\gamma\left(  \mathbf{r,\xi}\right)  =$ $D^{1}\left(  \mathbf{r,\xi;r,\xi
}\right)  $ \cite{weiner-trickey} and local external potentials of the
position variable $\mathbf{r}$ and spin space variable $\mathbf{\xi}$ that
describe the interaction with both external electric and magnetic fields.

For a system that contains on average $N$-electrons the space spin density
must satisfy the condition
\begin{equation}
\int\gamma\left(  \mathbf{r,\xi}\right)  d\mathbf{r}d\mathbf{\xi
}=N,\label{Number}%
\end{equation}
which does not restrict the density to be the contraction of a $N$-electron
state and one could consider incoherent and coherent mixtures of sates
describing different number of electrons. However in this article we focuse on
$N$-electron states and define convex sets, $\left\{  \left[  \gamma\right]
_{N}\right\}  ,$ of $N$-electron state density operators by
\begin{equation}
\left[  \gamma\right]  _{N}=\left\{  \left.  D\right\vert \Xi_{1}\left(
D\right)  =\gamma;D\in S_{N}\subset\mathcal{B}_{1}\left(  \mathcal{H}%
^{N}\right)  \right\}
\end{equation}
(see Appendix \ref{Denmat} for the definition of the convex set $S_{N}$ and
Appendix \ref{Operators} for the definition of $\mathcal{B}_{1}\left(
\mathcal{H}^{N}\right)  $). The constraining functional $\Xi_{1}$ can be
written in terms of the set of density operators $\left\{  \left.  \Phi\left(
\mathbf{x}\right)  ^{\dagger}\Phi\left(  \mathbf{x}\right)  \right\vert
\mathbf{x\in}\mathbb{X}\right\}  ,$ (see Appendix for the definition of the
field operator $\Phi\left(  \mathbf{x}\right)  $), as
\begin{equation}
\Xi_{1}\left(  D\right)  \left(  \mathbf{x}\right)  =\operatorname{Tr}\left\{
\Phi\left(  \mathbf{x}\right)  ^{\dagger}\Phi\left(  \mathbf{x}\right)
D\right\}  =\gamma\left(  \mathbf{x}\right)  .
\end{equation}
As $D$ is a positive operator it can always be expressed as the product
\begin{equation}
D=QQ^{\dagger},
\end{equation}
where $\left(  \text{Appendix \ref{Operators}}\right)  $
\begin{equation}
Q\in\mathcal{B}_{2}\left(  \mathcal{H}^{N}\right)  ,
\end{equation}
and $\mathcal{B}_{2}\left(  \mathcal{H}^{N}\right)  $ is the Hilbert space of
Hilbert-Schmidt operators acting on $N$-electron space $\mathcal{H}^{N}$.
However this decomposition is not unique, as in particular
\begin{equation}
D=QUU^{\dagger}Q^{\dagger}%
\end{equation}
for any unitary operator $U\in\mathcal{U}\left(  \mathcal{H}^{N}\right)  $.
One can remove this unitary group redundancy by choosing the operator $Q$ to
be self adjoint, leading to
\begin{equation}
D=A^{2},
\end{equation}
where
\begin{equation}
A\in\mathcal{B}_{2}\left(  \mathcal{H}^{N}\right)  \text{ and }A=A^{\dagger}.
\end{equation}
The operator $A$ is the square root of $D$, which again is not unique, as can
be observed from the spectral expansion
\begin{equation}
A=D^{\frac{1}{2}}=\sum_{1\leq j\leq r_{N}}\pm\sqrt{\alpha_{j}}\left\vert
\Upsilon_{j}\right\rangle \left\langle \Upsilon_{j}\right\vert ,\label{specA}%
\end{equation}
where
\begin{equation}
D=\sum_{1\leq j\leq r_{N}}\alpha_{j}\left\vert \Upsilon_{j}\right\rangle
\left\langle \Upsilon_{j}\right\vert ;\,\alpha_{j}\geq0,1\leq j\leq r_{N}%
\end{equation}
and $\left\{  \left.  \left\vert \Upsilon_{j}\right\rangle \right\vert 1\leq
j\leq r_{N}\right\}  $ is a complete orthornormal set of eigenvectors of $D$.
One could restrict $A$ to be positive i.e. use the positive square roots in
eq. $\left(  \ref{specA}\right)  ,$ but that would be too restrictive for our
purposes, as the set of positive Hilbert Schmidt operators forms a cone
\cite{cone} but do not form a linear space. However, the multi valuedness of
the square root function does not create a problem as the different roots are
separated from each other and could only form \emph{isolated} local minimizers
of the total energy functional. The domain of variation of the $N$-electron
energy functional can thus be taken to be the \emph{real} Hilbert space
$\mathcal{B}_{2sa}\left(  \mathcal{H}^{N}\right)  ,$ of dimension $r_{N}^{2},$
which has orthonormal bases
\begin{align}
\Lambda_{jk} &  =\frac{1}{\sqrt{2}}\left\{  \left\vert \Psi_{j}\right\rangle
\left\langle \Psi_{k}\right\vert +\left\vert \Psi_{k}\right\rangle
\left\langle \Psi_{j}\right\vert \right\}  ;\quad1\leq j<k\leq r_{N}%
\nonumber\\
\Lambda_{jj} &  =\left\vert \Psi_{j}\right\rangle \left\langle \Psi
_{j}\right\vert ;\quad1\leq j\leq r_{N}\nonumber\\
\Lambda_{jk} &  =\frac{i}{\sqrt{2}}\left\{  \left\vert \Psi_{j}\right\rangle
\left\langle \Psi_{k}\right\vert -\left\vert \Psi_{k}\right\rangle
\left\langle \Psi_{j}\right\vert \right\}  ;\quad1\leq j>k\leq r_{N},
\end{align}
where $\left\{  \left.  \left\vert \Psi_{j}\right\rangle \right\vert 1\leq
j\leq r_{N}\right\}  $ is any conb for $\mathcal{H}^{N},$ of course all these
spaces are infinite dimensional if $r_{N}=\infty.$

\subsubsection{The Density Constraint}

The Levy-Leib construction is based on defining the value of the density
functional as the minimum of a energy variation over states constrained to
have a fixed density. The density constraint can be formulated in terms of a
quadratic functional $\Xi_{2}$ based on the linear functional $\Xi_{1},$
where
\begin{equation}
\Xi_{2}:\mathcal{B}_{2sa}\left(  \mathcal{H}^{N}\right)  \rightarrow
L_{1+}\left(  \mathbb{X}\right)  \subset L_{1}\left(  \mathbb{X}\right)  ,
\end{equation}
$L_{1+}\left(  \mathbb{X}\right)  $ is the positive cone of $L_{1}\left(
\mathbb{X}\right)  ,$ i.e. the set of positive functions$,$ it is given by
\begin{equation}
\Xi_{2}\left(  A\right)  =\operatorname{Tr}\left\{  A\Phi^{\dagger}\Phi
A\right\}  =\left(  A\left\vert \Phi^{\dagger}\Phi A\right.  \right)  =\gamma,
\end{equation}
where the real symmetric inner product $\left(  \cdot\left\vert \cdot\right.
\right)  $ in $\mathcal{B}_{2sa}\left(  \mathcal{H}^{N}\right)  $ is defined
as
\begin{equation}
\left(  X\left\vert Y\right.  \right)  =\operatorname{Tr}\left\{  XY\right\}
.
\end{equation}
In terms of the linear functional $\Xi_{1}$ it is given by
\begin{equation}
\Xi_{1}\left(  X\right)  =\operatorname{Tr}\left\{  \Phi^{\dagger}\Phi
X\right\}  =\gamma\,
\end{equation}
and one must impose the further constraint $X\geq0.$ If the differentials of
the component constraints $\left\{  \Xi_{2}\left(  A\right)  \left(
\mathbf{x}\right)  \right\}  $, at a point $X$ form a basis (see Appendix
basis \ref{Bases}) for $L_{1}\left(  \mathbb{X}\right)  $ then this point is
called \emph{regular}. The components are
\begin{align}
\left[  D\Xi_{2}\left(  A\right)  \left(  \mathbf{x}\right)  \right]  \left(
\cdot\right)   &  =\operatorname{Tr}\left\{  A\Phi\left(  \mathbf{x}\right)
^{\dagger}\Phi\left(  \mathbf{x}\right)  \cdot\right\}  =\left(  A\left\vert
\Phi\left(  \mathbf{x}\right)  ^{\dagger}\Phi\left(  \mathbf{x}\right)
\right.  \cdot\right)  \\
\left[  D\Xi_{1}\left(  X\right)  \left(  \mathbf{x}\right)  \right]  \left(
\cdot\right)   &  =\left[  \Xi_{1}\left(  \mathbf{x}\right)  \right]  \left(
\cdot\right)  =\operatorname{Tr}\left\{  \Phi\left(  \mathbf{x}\right)
^{\dagger}\Phi\left(  \mathbf{x}\right)  \cdot\right\}  .
\end{align}


It will be convenient to use multiplication super operators
\begin{equation}
\widehat{T}:\mathcal{B}_{2sa}\left(  \mathcal{H}^{N}\right)  \rightarrow
\mathcal{B}_{2sa}\left(  \mathcal{H}^{N}\right)  ,
\end{equation}
which could be unbounded, and are defined by the quadratic form
\cite{weiner05}
\begin{equation}
\left(  X\left\vert \widehat{T}Y\right.  \right)  =\operatorname{Tr}\left\{
XTY\right\}  \,\forall X\in\mathcal{B}_{2sa}\left(  \mathcal{H}^{N}\right)
,\,Y\in\operatorname{Dom}\left\{  \widehat{T}\right\}  ,
\end{equation}
where
\[
\operatorname{Dom}\left\{  \widehat{T}\right\}  =\left\{  \left.  Y\right\vert
\operatorname{Tr}\left\{  XTY\right\}  <\infty\,\forall X\in\mathcal{B}%
_{2sa}\left(  \mathcal{H}^{N}\right)  \right\}  .
\]
This gives
\begin{equation}
\widehat{T}\left\vert Y\right)  =\left\vert \left[  TY\right]  \right)  ,
\end{equation}
where
\begin{equation}
T:\operatorname{Dom}\left\{  T\right\}  \subseteq\mathcal{H}^{N}%
\rightarrow\mathcal{H}^{N}%
\end{equation}
and $\left\vert \left[  TY\right]  \right)  $ is the projection of $TY$ onto
$\mathcal{B}_{2sa}\left(  \mathcal{H}^{N}\right)  $ i.e. the self adjoint part
of $TY,$ given by
\[
\left[  TY\right]  =\frac{1}{2}\left\{  TY+YT\right\}  =\frac{1}{2}\left[
T,Y\right]  _{+}.
\]
One can observe that
\begin{equation}
\operatorname{Tr}\left\{  XTY\right\}  =\operatorname{Tr}\left\{  X\left[
TY\right]  \right\}  ,
\end{equation}
as $X\in\mathcal{B}_{2sa}\left(  \mathcal{H}^{N}\right)  $ and is orthogonal
to the subspace of anti self adjoint operators.

\subsubsection{Hohenberg-Kohn Density Functional Construction}

Even though the explicit form of the Hohenberg-Kohn functional $\mathcal{F}%
_{HK}\left(  \gamma\right)  $ is not known one can define it implicitly
\cite{LevyLieb} as
\begin{equation}
\mathcal{F}_{HK}\left(  \gamma\right)  =\operatorname{Tr}\left\{  \left(
K+V_{2}\right)  D_{\ast}\left(  \gamma\right)  \right\}  =E_{U}\left(
D_{\ast}\left(  \gamma\right)  \right)
\end{equation}
in terms of the minima, $E_{U}\left(  D_{\ast}\left(  \gamma\right)  \right)
,$ of the universal energy functional $E_{U}\left(  D\right)  ,$ that are
solutions of the constrained linear optimization
\begin{equation}
\min_{D\in\mathcal{B}_{1sa}\left(  \mathcal{H}^{N}\right)  }\operatorname{Tr}%
\left\{  \left(  K+V_{2}\right)  D\right\}  \equiv\min_{D\in\mathcal{B}%
_{1sa}\left(  \mathcal{H}^{N}\right)  }\text{ }E_{U}\left(  D\right)
\label{HK}%
\end{equation}
subject to
\begin{align}
\Xi_{1}\left(  D\right)   &  =\gamma,\nonumber\\
\operatorname{Tr}D  &  =1,\nonumber\\
D  &  \geq0, \label{HKC}%
\end{align}
which is a linear programing problem, or by the constrained quadratic
optimization
\begin{equation}
\min_{A\in\mathcal{B}_{2sa}\left(  \mathcal{H}^{N}\right)  }\operatorname{Tr}%
\left\{  A\overset{}{\left(  K+V_{2}\right)  A}\right\}  =\left(  A_{\ast
}\left(  \gamma\right)  \left\vert \widehat{\left(  K+V_{2}\right)  }A_{\ast
}\left(  \gamma\right)  \right.  \right)  =E_{U}\left(  A_{\ast}\left(
\gamma\right)  \right)  \label{HKQ}%
\end{equation}
subject to
\begin{align}
&  \left(  A\left\vert \widehat{\Phi^{\dagger}\Phi}A\right.  \right)
=\gamma,\nonumber\\
&  \Xi_{2}\left(  A\right)  =\left(  A\left\vert \widehat{\Phi^{\dagger}\Phi
}A\right.  \right)  =\gamma,
\end{align}
where
\begin{equation}
D_{\ast}\left(  \gamma\right)  =A_{\ast}\left(  \gamma\right)  ^{2}.
\end{equation}
We will focus on the quadratic formulation, which is a convex constrained
optimization problem and therefore a local minimum is a global minimum, (which
may be degenerate). The Lagrangian associated with this problem eq. $\left(
\ref{HK}\right)  $ is
\begin{equation}
\mathcal{L}\left(  A,\lambda,\gamma\right)  =E_{U}\left(  A\right)
-\left\langle \left.  \lambda\right\vert \left\{  \Xi_{2}\left(  A\right)
-\gamma\right\}  \right\rangle , \label{Lar}%
\end{equation}
where the universal energy is
\begin{equation}
E_{U}\left(  A\right)  =\operatorname{Tr}\left\{  \left(  K+V_{2}\right)
A^{2}\right\}  =\left(  A\left\vert \widehat{\left(  K+V_{2}\right)
}A\right.  \right)  ,
\end{equation}
the Lagrange multiplier function is
\begin{equation}
\left\langle \left.  \lambda\right\vert \left\{  \Xi_{2}\left(  A\right)
-\gamma\right\}  \right\rangle =\int\lambda\left(  \mathbf{x}\right)  \left\{
\Xi_{2}\left(  A\right)  \left(  \mathbf{x}\right)  -\gamma\left(
\mathbf{x}\right)  \right\}  d\mathbf{x,\,}\lambda\in L_{\infty}\left(
\mathbb{X}\right)  ,\gamma\in L_{1}\left(  \mathbb{X}\right)  ,
\end{equation}
where $\left\langle \cdot|\cdot\right\rangle $ is the scalar product between
the dual spaces $\left\langle L_{\infty}\left(  \mathbb{X}\right)
,L_{1}\left(  \mathbb{X}\right)  \right\rangle $ i.e.
\begin{equation}
\left\langle \cdot|\cdot\right\rangle :L_{\infty}\left(  \mathbb{X}\right)
\times L_{1}\left(  \mathbb{X}\right)  \rightarrow\mathbb{R}%
\end{equation}
and the control density $\gamma$ is chosen to belong to the convex set
$\mathcal{P}_{N}$ given by
\begin{equation}
\mathcal{P}_{N}=\left\{  f\left\vert f\in L_{1}\left(  \mathbb{X}\right)
;\,f\left(  \mathbf{x}\right)  \geq0;\int f\left(  \mathbf{x}\right)
d\mathbf{x}=N\right.  \right\}  .
\end{equation}
This choice of the density $\gamma$ coupled with insisting that $A\in
\mathcal{B}_{2sa}\left(  \mathcal{H}^{N}\right)  $ ensures that $A$ is
normalized (Appendix \ref{Normalization}$),$ i.e.
\begin{equation}
\left(  A\left\vert A\right.  \right)  =1,
\end{equation}
thus the normalization constraint is not an active constraint at such
densities and the Lagrange parameter associated with it would be zero.

The constraint term in the Lagrangian eq. $\left(  \ref{Lar}\right)  $ can be
developed to give
\begin{equation}
\left\langle \lambda\left\vert \gamma\right.  \right\rangle -\left(
A\left\vert \widehat{\Omega\left(  \lambda\right)  }\right.  A\right)
=\left(  A\left\vert \left\{  \left\langle \lambda\left\vert \gamma\right.
\right\rangle \hat{I}-\widehat{\Omega\left(  \lambda\right)  }\right\}
\right.  A\right)  ,
\end{equation}
where $\Omega\left(  \lambda\right)  $ is a local potential defined as
\begin{equation}
\Omega\left(  \lambda\right)  =\int\lambda\left(  \mathbf{x}\right)
\mathbf{\Phi}\left(  \mathbf{x}\right)  ^{\dagger}\mathbf{\Phi}\left(
\mathbf{x}\right)  d\mathbf{x}%
\end{equation}
and
\begin{equation}
\left(  A\left\vert \widehat{\Omega\left(  \lambda\right)  }\right.  A\right)
=\int\lambda\left(  \mathbf{x}\right)  \Xi_{2}\left(  A\right)  \left(
\mathbf{x}\right)  d\mathbf{x.}%
\end{equation}


\subsubsection{Optimization Conditions}

Solutions $A_{\ast}$ of the constrained quadratic optimization eq. $\left(
\ref{HKQ}\right)  $ correspond to unconstrained local saddle points $\left(
A_{\ast},\lambda_{\ast}\right)  $ \cite{KKT} of the Lagrangian $\mathcal{L}%
\left(  A_{\ast},\lambda_{\ast},\gamma\right)  $
\begin{equation}
\mathcal{L}\left(  A_{\ast},\lambda_{\ast},\gamma\right)  =\max_{\lambda}%
\min_{A}\mathcal{L}\left(  A,\lambda,\gamma\right)
\end{equation}
that satisfy the necessary condition
\begin{equation}
D\mathcal{L}\left(  A_{\ast},\lambda_{\ast},\gamma\right)  =0,
\end{equation}
where the linear map
\begin{equation}
D\mathcal{L}\left(  A,\lambda,\gamma\right)  :\mathcal{B}_{2sa}\left(
\mathcal{H}^{N}\right)  \times L_{\infty}\left(  \mathbb{X}\right)
\rightarrow\mathbb{R}%
\end{equation}
is given by%
\begin{equation}
D\mathcal{L}\left(  A,\lambda,\gamma\right)  \left(  h,\varphi\right)
=\left\langle \varphi\left\vert \gamma\right.  \right\rangle -\left(
A\left\vert \widehat{\Omega\left(  \varphi\right)  }\right.  A\right)
+2\left(  h\left\vert \left\{  \widehat{K+V_{2}-\Omega\left(  \lambda\right)
}\right\}  \right.  A\right)  . \label{LagDiff}%
\end{equation}
We have shown \cite{weiner05} that in general $\lambda_{\ast}$ is unique, but
there may be degenerate $A_{\ast}$'s. The differential eq. $\left(
\ref{LagDiff}\right)  $ can be partitioned into $\left(  D_{A}\mathcal{L}%
\left(  A,\lambda,\gamma\right)  ,D_{\lambda}\mathcal{L}\left(  A,\lambda
,\gamma\right)  \right)  $, where
\begin{align}
D_{A}\mathcal{L}\left(  A,\lambda,\gamma\right)   &  =2\left(  \cdot\left\vert
\widehat{\left\{  K+V_{2}-\Omega\left(  \lambda\right)  \right\}  }\right.
A\right) \nonumber\\
D_{\lambda}\mathcal{L}\left(  A,\lambda,\gamma\right)   &  =\left\langle
\cdot\left\vert \gamma\right.  \right\rangle -\left(  A\left\vert
\Omega\left(  \cdot\right)  \right.  A\right)  ,
\end{align}
thus giving the stationary point conditions
\begin{subequations}
\begin{align}
\left(  h\left\vert \widehat{\left\{  K+V_{2}-\Omega\left(  \lambda_{\ast
}\right)  \right\}  }\right.  A_{\ast}\right)   &  =0\,\forall h\in
\mathcal{B}_{2sa}\left(  \mathcal{H}^{N}\right) \label{stata}\\
\left\langle \varphi\left\vert \gamma\right.  \right\rangle -\left(  A_{\ast
}\left\vert \widehat{\Omega\left(  \varphi\right)  }\right.  A_{\ast}\right)
&  =0\,\forall\varphi\in L_{\infty}\left(  \mathbb{X}\right)  . \label{statb}%
\end{align}
If eq. $\left(  \ref{stata}\right)  $ and eq. $\left(  \ref{statb}\right)  $
are true for all operators $h$ belonging to $\mathcal{B}_{2sa}\left(
\mathcal{H}^{N}\right)  $ and all functions $\varphi$ belonging to $L_{\infty
}\left(  \mathbb{X}\right)  $ these conditions are equivalent to
\end{subequations}
\begin{equation}
\widehat{\left\{  K+V_{2}-\Omega\left(  \lambda_{\ast}\right)  \right\}
}\left\vert A_{\ast}\right)  =0 \label{Stateig}%
\end{equation}
and constraining condition
\begin{equation}
\left(  A_{\ast}\left\vert \widehat{\Phi^{\dagger}\Phi}\right.  A_{\ast
}\right)  =\gamma.
\end{equation}
(For unbounded operators $K+V_{2}-\Omega\left(  \lambda_{\ast}\right)  $ one
must be more precise, but here we will only consider the case when an
eigenvector $\left\vert A_{\ast}\right)  $ exists and thus eq. $\left(
\ref{stata}\right)  $ is true for all $h\in\mathcal{B}_{2sa}\left(
\mathcal{H}^{N}\right)  $ and $\gamma$ is V-representable$,$ more general
cases are discussed in \cite{weiner2006}). The optimal value of the Lagrangian
satisfies the following identities:
\begin{align}
\mathcal{L}\left(  A_{\ast},\lambda_{\ast},\gamma\right)   &  =\left(
A_{\ast}\left(  \gamma\right)  \left\vert \left\{  \widehat{\left\{
K+V_{2}-\Omega\left(  \lambda_{\ast}\right)  \right\}  }+\left\langle \left.
\lambda_{\ast}\left(  \gamma\right)  \right\vert \gamma\right\rangle \hat
{I}\right\}  \right.  A_{\ast}\left(  \gamma\right)  \right)  =\left\langle
\left.  \lambda_{\ast}\left(  \gamma\right)  \right\vert \gamma\right\rangle
\nonumber\\
&  =\left(  A_{\ast}\left(  \gamma\right)  \left\vert \widehat{K+V_{2}%
}\right.  A_{\ast}\left(  \gamma\right)  \right)  =E_{U}\left(  A_{\ast
}\left(  \gamma\right)  \right)  .
\end{align}
This shows that the optimal values of the universal energy functional, the
value of the Lagrangian and the value of the HK energy functional are all
equal at the density $\gamma$ i.e.
\begin{equation}
\mathcal{F}_{HK}\left(  \gamma\right)  =E_{U}\left(  A_{\ast}\left(
\gamma\right)  \right)  =\left\langle \left.  \lambda_{\ast}\left(
\gamma\right)  \right\vert \gamma\right\rangle .
\end{equation}
Moreover $\gamma$ is the density of the ground $D_{\ast}\left(  \gamma\right)
=A_{\ast}\left(  \gamma\right)  ^{2}$ of the Hamiltonian $K+V_{2}%
-\Omega\left(  \lambda_{\ast}\left(  \gamma\right)  \right)  $ i.e. of an
electronic system in an external potential $\lambda_{\ast}\left(
\gamma\right)  $.

\subsubsection{Sensitivity Analysis}

Sensitivity analysis \cite{sensitivity} examines how the optimal solution
changes with variation of the value of the constraining functions, which can
be thought of as control parameters. In this case the $N$-electron state
energy with variations of $\gamma.$ If the optimal points $\left\{  A_{\ast
},\lambda_{\ast},\gamma\right\}  $ are regular for a range of $\gamma^{\prime
}s$ then $\left\{  A_{\ast},\lambda_{\ast}\right\}  $ are in fact functions of
$\gamma,$ i.e.
\begin{align}
A_{\ast}  &  =A_{\ast}\left(  \gamma\right) \nonumber\\
\lambda_{\ast}  &  =\lambda_{\ast}\left(  \gamma\right)  ,
\end{align}
and are single valued. The sensitivity of the optimal $N$-electron energy to
changes in the density $\gamma$ is then given by
\begin{equation}
\frac{\delta E_{U}\left(  A_{\ast}\left(  \gamma\right)  \right)  }%
{\delta\gamma}=\lambda_{\ast}\left(  \gamma\right)  .
\end{equation}
By definition
\begin{equation}
\mathcal{F}_{HK}\left(  \gamma\right)  =E_{U}\left(  A_{\ast}\left(
\gamma\right)  \right)  ,
\end{equation}
which shows
\begin{equation}
\frac{\delta\mathcal{F}_{HK}\left(  \gamma\right)  }{\delta\gamma}%
=\lambda_{\ast}\left(  \gamma\right)  .
\end{equation}


\subsubsection{Decomposition of the Hohenburg-Kohn Functional}

If a minimizer $D_{\ast}\left(  \gamma\right)  $ exists and is unique and then
one can define kinetic energy, coulomb and exchange-correlation functionals by
eq. $\left(  \ref{HK}\right)  $ as
\begin{align}
\mathcal{F}_{K}\left(  \gamma\right)   &  =\operatorname{Tr}\left\{  KD_{\ast
}\left(  \gamma\right)  \right\}  ,\nonumber\\
\mathcal{F}_{C}\left(  \gamma\right)   &  =\int\frac{\gamma\left(
\mathbf{r,\xi}\right)  \gamma\left(  \mathbf{r}^{\prime}\mathbf{,\xi}\right)
}{\left\Vert \mathbf{r-r}^{\prime}\right\Vert }d\mathbf{r}d\mathbf{r}^{\prime
}d\mathbf{\xi},\nonumber\\
\mathcal{F}_{XC}\left(  \gamma\right)   &  =\operatorname{Tr}\left\{
V_{2}D_{\ast}\left(  \gamma\right)  \right\}  -\mathcal{F}_{C}\left(
\gamma\right)  ,
\end{align}
producing the decomposition
\begin{equation}
\mathcal{F}_{HK}\left(  \gamma\right)  =\mathcal{F}_{K}\left(  \gamma\right)
+\mathcal{F}_{C}\left(  \gamma\right)  +\mathcal{F}_{XC}\left(  \gamma\right)
.
\end{equation}
The concept of a unique exchange-correlation functional is central to the
theory, but the existence of non isolated degenerate $D_{\ast}\left(
\gamma\right)  ^{\prime}$s could give different $\mathcal{F}_{XC}\left(
\gamma\right)  ^{\prime}$s and $\mathcal{F}_{K}\left(  \gamma\right)
^{\prime}$s. However if degeneracies do occur one can still define an
exchange-correlation functional through the imposition of further constrained
minimizations \cite{weiner06}.

\subsubsection{Molecular System Energy}

The total energy functional of a system that contains $N$ electrons and
experiences an environment described by a \emph{local} potential $\mathrm{v}$
is
\begin{equation}
E\left(  \gamma\mathbf{,}\mathrm{v}\right)  =\mathcal{F}_{HK}\left(
\gamma\right)  +f_{\gamma}\left(  \mathrm{v}\right)  ,
\end{equation}
where the explicit effect of the environment is described by the functional
\begin{equation}
f_{\gamma}\left(  \mathrm{v}\right)  =\int\mathrm{v}\left(  \mathbf{r,\xi
}\right)  \gamma\left(  \mathbf{r,\xi}\right)  d\mathbf{r}d\mathbf{\xi.}%
\end{equation}
The ground state energy $E_{GS}\left(  \mathrm{v}\right)  $ and the
corresponding space-spin density $\gamma_{GS}$ $\left(  \mathrm{v}\right)  $
are then determined by
\begin{equation}
E_{GS}\left(  \mathrm{v}\right)  =E\left(  \gamma_{GS}\mathbf{,}%
\mathrm{v}\right)  =\min_{\mathbf{\gamma\in}\mathcal{D}}E\left(
\gamma\mathbf{,}\mathrm{v}\right)  , \label{GS}%
\end{equation}
where
\begin{equation}
\mathcal{D}=\left\{  \gamma\left\vert \int\gamma\left(  \mathbf{r,\xi}\right)
d\mathbf{r}d\mathbf{\xi}=N,\quad\gamma\left(  \mathbf{r,\xi}\right)
\geq0,\,\gamma\in L_{1}\left(  \mathbb{X}\right)  \right.  \right\}  .
\end{equation}
If $\gamma$ is any V-representable density one can deduce from eq. $\left(
\ref{Stateig}\right)  $ that
\begin{equation}
\widehat{\left\{  K+V_{2}+\mathrm{v}-\Omega\left(  \lambda_{\ast}\left(
\gamma\right)  \right)  \right\}  }\left\vert A_{\ast}\left(  \gamma\right)
\right)  =0 \label{vstat}%
\end{equation}
as $f_{\gamma}\left(  \mathrm{v}\right)  $ is just a constant for fixed
$\gamma$, which can be recast as
\begin{equation}
\widehat{\left\{  K+V_{2}-\Omega\left(  \lambda_{\ast}\left(  \gamma\right)
-\mathrm{v}\right)  \right\}  }\left\vert A_{\ast}\left(  \gamma\right)
\right)  =0,
\end{equation}
where
\begin{equation}
\Omega\left(  \lambda_{\ast}\left(  \gamma\right)  -\mathrm{v}\right)
=\int\left\{  \lambda_{\ast}\left(  \gamma\right)  -\mathrm{v}\right\}
\left(  \mathbf{x}\right)  \Phi\left(  \mathbf{x}\right)  ^{\dagger}%
\Phi\left(  \mathbf{x}\right)  d\mathbf{x.}%
\end{equation}


In order to find the ground state energy $E\left(  \ \gamma_{GS}%
\mathbf{,}\mathrm{v}\right)  $ one needs to solve a constrained optimization
problem eq. $\left(  \ref{GS}\right)  $,
\begin{equation}
\min_{\gamma}E\left(  \gamma\mathbf{,}\mathrm{v}\right)  ;\quad\gamma
\in\mathcal{D},
\end{equation}
which one can make unconstrained by setting
\begin{equation}
\gamma\left(  \mathbf{x}\right)  =\frac{Nf\left(  \mathbf{x}\right)  ^{2}%
}{\left\langle f|f\right\rangle }=\frac{N\left\langle f\left\vert \Phi\left(
\mathbf{x}\right)  ^{\dagger}\Phi\left(  \mathbf{x}\right)  f\right.
\right\rangle }{\left\langle f|f\right\rangle },\,f\in L_{2}\left(
\mathbb{X}\right)  .
\end{equation}
The stationary points of $E\left(  f\mathbf{,}\mathrm{v}\right)  $ with
respect to variations of $f$ satisfy
\begin{equation}
\frac{\delta E\left(  f\mathbf{,}\mathrm{v}\right)  }{\delta f\left(
\mathbf{x}\right)  }=\int\frac{\delta E\left(  f\mathbf{,}\mathrm{v}\right)
}{\delta\gamma\left(  \mathbf{x}^{\prime}\right)  }\frac{\delta\gamma\left(
\mathbf{x}^{\prime}\right)  }{\delta f\left(  \mathbf{x}\right)  }%
d\mathbf{x}^{\prime}=\int\left\{  \frac{\delta\mathcal{F}_{HK}\left(
\gamma\right)  }{\delta\gamma\left(  \mathbf{x}^{\prime}\right)  }%
+\mathrm{v}\left(  \mathbf{x}^{\prime}\right)  \right\}  \frac{\delta
\gamma\left(  \mathbf{x}^{\prime}\right)  }{\delta f\left(  \mathbf{x}\right)
}=0\text{ for almost all }\mathbf{x} \label{Efderiv}%
\end{equation}
Thus one obtains the well known result
\begin{equation}
\frac{\delta\mathcal{F}_{HK}\left(  \gamma\right)  }{\delta\gamma\left(
\mathbf{x}^{\prime}\right)  }\overset{a.e}{=}-\mathrm{v}\left(  \mathbf{x}%
^{\prime}\right)  .
\end{equation}


\subsection{General Spin Orbital Variables \label{GSO}}

There are two standard approaches to formulating the total energy in terms of
GSO's and occupation numbers, one is to consider functionals of FORDO's, which
leads to DMFT, the other is to factor the dependence through the density,
which leads to KS-DFT and is based on the HK functional. However as the
explicit forms of these functionals are not known one must rely on the
implicit definitions via constrained optimizations, though in practice
approximate explicit forms are utilized. We have recently \cite{WeinerOrtiz05}
introduced a third approach which is defined on states generated by the
submanifold of $2N$-electron AGP states. The total energy functional, which is
explicitly given, \emph{is not} a density matrix functional though it
\emph{does} express the total energy in terms of GSO's and occupation numbers.

\subsubsection{Correspondence between Orbitals and Densities}

The density $\gamma_{GS}$ $\left(  \mathrm{v}\right)  $ can be always be
expressed in terms of the Natural General Spin Orbitals (NGSO's) $\left\{
\tilde{\psi}_{j}\right\}  $ and occupation numbers $\left\{  N_{j}\right\}  $
of the FORDO $D_{\ast}^{1}\left(  \gamma\right)  $ (Appendix \ref{Denmat}) as%

\begin{equation}
\left[  \gamma_{GS}\left(  \mathrm{v}\right)  \right]  \left(  \mathbf{x}%
\right)  =\left[  D_{\ast}^{1}\left(  \gamma_{GS}\left(  \mathrm{v}\right)
\right)  \right]  \left(  \mathbf{x,x}\right)  =\sum_{1\leq j\leq r_{1}}%
N_{j}\left\vert \tilde{\psi}_{j}\left(  \mathbf{x}\right)  \right\vert ^{2},
\label{NGSOexpan}%
\end{equation}
but many other sets of GSO's also produce the same density, in particular
Harimanns \cite{Harimann} theorem ensures that there are always expansions of
the form
\begin{equation}
\left[  \gamma_{GS}\left(  \mathrm{v}\right)  \right]  \left(  \mathbf{x}%
\right)  =\left[  D_{\ast}^{1}\left(  \gamma_{GS}\left(  \mathrm{v}\right)
\right)  \right]  \left(  \mathbf{x,x}\right)  =\sum_{1\leq j\leq N}\left\vert
\varphi_{j}\left(  \mathbf{x}\right)  \right\vert ^{2} \label{Genexp}%
\end{equation}
in terms of orthonormal square summable functions $\left\{  \left.
\varphi_{j}\right\vert 1\leq j\leq N\right\}  $.

In general a density $\gamma$ can be expressed in two ways, one in terms of
unnormalized orthogonal sets of GSO's $\mathbf{\varphi,}$ treating the
occupation numbers $\mathbf{N}$ as control parameters, the other in terms of
orthonormal sets $\mathbf{\tilde{\varphi}}$ and occupation numbers
$\mathbf{N,}$ viewed as constrained independent variables. This is made
explicit through the introduction of two functionals $\Gamma_{1}$ and
$\Gamma_{2}.$ The first functional $\Gamma_{1}$ is defined by
\begin{equation}
\Gamma_{1}:\mathfrak{S}\left(  \mathbf{N}\left(  N\right)  \right)
\subset\mathcal{H}^{1}\rightarrow L_{1+}\left(  \mathbb{X}\right)  \subset
L_{1}\left(  \mathbb{X}\right)  ,
\end{equation}
where%
\begin{align}
\gamma &  =\Gamma_{1}\left(  \mathbf{\varphi}\right)  ,\nonumber\\
\left[  \Gamma_{1}\left(  \mathbf{\varphi}\right)  \right]  \left(
\mathbf{x}\right)   &  =\sum_{1\leq j\leq r_{1}}\left\vert \varphi_{j}\left(
\mathbf{x}\right)  \right\vert ^{2}=\gamma\left(  \mathbf{x}\right)  .
\label{Gamma}%
\end{align}
and the domain $\mathfrak{S}\left(  \mathbf{N}\left(  N\right)  \right)  $ is
given by
\begin{equation}
\mathfrak{S}_{1}\left(  \mathbf{N}\left(  N\right)  \right)  =\left\{
\mathbf{\varphi\in}\mathfrak{S}_{1}\left(  \mathbf{N}\right)  \left\vert
\sum_{1\leq j\leq r_{1}}N_{j}=N\right.  \right\}  ,
\end{equation}
where
\begin{equation}
\mathfrak{S}_{1}\left(  \mathbf{N}\right)  =\left\{  \mathbf{\varphi
}\left\vert \left\langle \left.  \varphi_{j}\right\vert \varphi_{k}%
\right\rangle =N_{j}\delta_{jk};\,0\leq N_{j}\leq1,1\leq j,k\leq r_{1}\right.
\right\}
\end{equation}
and as many as $r_{1}-N$ of the component functions, $\varphi_{j},$ may be
identically zero. The second functional $\Gamma_{2}$ is defined by
\begin{equation}
\Gamma_{2}:\mathfrak{S}_{2}\left(  N\right)  \times\mathfrak{N}\left(
N\right)  \subset\mathcal{H}^{1}\times\mathbb{R}^{r_{1}}\rightarrow
L_{1+}\left(  \mathbb{X}\right)  \subset L_{1}\left(  \mathbb{X}\right)  ,
\end{equation}
where
\begin{equation}
\Gamma_{2}\left(  \mathbf{\tilde{\varphi},N}\right)  =\sum_{1\leq j\leq r_{1}%
}\left\vert \tilde{\varphi}_{j}\left(  \mathbf{x}\right)  \right\vert
^{2}N_{j}=\gamma\left(  \mathbf{x}\right)  , \label{GammaN}%
\end{equation}
$\mathfrak{S}_{2}\left(  N\right)  $ is composed of sets of orthonormalized
GSO's
\begin{equation}
\mathfrak{S}_{2}=\left\{  \mathbf{\tilde{\varphi}}\left\vert \left\langle
\left.  \tilde{\varphi}_{j}\right\vert \tilde{\varphi}_{k}\right\rangle
=\delta_{jk};\,1\leq j,k\leq r_{1}\right.  \right\}  ,
\end{equation}
and the set of real vectors $\mathfrak{N}\left(  N\right)  \subset
\mathbb{R}^{r_{1}}$ given by
\begin{equation}
\mathfrak{N}\left(  N\right)  =\left\{  \mathbf{N}\left\vert 0\leq N_{j}%
\leq1,1\leq j\leq r_{1};\sum_{1\leq j\leq r_{1}}N_{j}=N\right.  \right\}  .
\end{equation}
Clearly
\begin{equation}
\gamma=\Gamma_{1}\left(  \mathbf{\varphi}\right)  =\Gamma_{2}\left(
\mathbf{\tilde{\varphi},N}\right)  ,
\end{equation}
where $\mathbf{\varphi}$ and $\left\{  \mathbf{\tilde{\varphi},N}\right\}  $
determine the same FORDO $D^{1}$ i.e.
\begin{equation}
D^{1}\left(  \mathbf{\varphi}\right)  =D^{1}\left(  \mathbf{\tilde{\varphi}%
,N}\right)  =\sum_{1\leq j\leq r_{1}}\left\vert \varphi_{j}\right\rangle
\left\langle \varphi_{j}\right\vert =\sum_{1\leq j\leq r_{1}}\left\vert
\tilde{\varphi}_{j}\right\rangle \left\langle \tilde{\varphi}_{j}\right\vert
N_{j}.
\end{equation}
The optimization wrt $\gamma$ can thus be expressed in terms of
$\mathbf{\varphi}$ via $\Gamma_{1}$ or $\left\{  \mathbf{\tilde{\varphi}%
,N}\right\}  $ via $\Gamma_{2}.$ We will base the following analysis on
$\Gamma_{1},$ but describe and compare in Appendix \ref{GSO&Occ} the analogous
analysis in terms of $\Gamma_{2}.$

The functional $\Gamma_{1}$ is not one to one and is constant on equivalence
classes, $\mathfrak{S}\left(  \mathbf{N}\left(  N\right)  \mathbf{,\gamma
}\right)  ,$ of sets of GSO's $\mathbf{\varphi}$ corresponding to the
equivalence classes $\left[  \gamma\right]  _{1}$ of FORDO's, that produce the
same density. Hence an inverse functional $\Gamma_{1}^{-1}$ does not exist. We
will, however, discuss later a procedure that essentially provides a one to
one relationship between densities and GSO's and thus leads to a generalized inverse.

\subsubsection{GSO functional}

One can express the total energy via the orbital functional $\mathcal{F}_{O}$,
factored through the HK energy functional, as
\begin{equation}
\mathcal{F}_{O}\left(  \mathbf{\varphi}\right)  =\mathcal{F}_{HK}\left(
\Gamma_{1}\left(  \mathbf{\varphi}\right)  \right)  =\mathcal{F}_{HK}\left(
\gamma\right)  . \label{FO}%
\end{equation}


The functional $\mathcal{F}_{O}$ can be decomposed as
\begin{equation}
\mathcal{F}_{O}\left(  \mathbf{\varphi}\right)  =\mathcal{F}_{T}\left(
\mathbf{\varphi}\right)  +\mathcal{F}_{C}\left(  \mathbf{\varphi}\right)
+\mathcal{F}_{OXC}\left(  \mathbf{\varphi}\right)  , \label{KS1}%
\end{equation}
where the orbital kinetic energy functional $\mathcal{F}_{T}$
\begin{equation}
\mathcal{F}_{T}\left(  \mathbf{\varphi}\right)  =\sum_{1\leq j\leq r_{1}%
}\left\langle \varphi_{j}\left\vert K\varphi_{j}\right.  \right\rangle ,
\end{equation}
the orbital exchange-corellation functional $\mathcal{F}_{OXC}$ is defined as
\begin{equation}
\mathcal{F}_{OXC}\left(  \mathbf{\varphi}\right)  =\mathcal{F}_{XC}\left(
\Gamma_{1}\left(  \mathbf{\varphi}\right)  \right)  +\mathcal{F}_{K}\left(
\Gamma_{1}\left(  \mathbf{\varphi}\right)  \right)  -\mathcal{F}_{T}\left(
\mathbf{\varphi}\right)  , \label{KS2}%
\end{equation}
and one recalls the definitions
\begin{equation}
\mathcal{F}_{XC}\left(  \Gamma_{1}\left(  \mathbf{\varphi}\right)  \right)
=\operatorname{Tr}\left\{  V_{2}D_{\ast}\left(  \Gamma_{1}\left(
\mathbf{\varphi}\right)  \right)  \right\}  -\mathcal{F}_{C}\left(  \Gamma
_{1}\left(  \mathbf{\varphi}\right)  \right)  \label{KS3}%
\end{equation}
and
\begin{equation}
\mathcal{F}_{K}\left(  \Gamma_{1}\left(  \mathbf{\varphi}\right)  \right)
=\operatorname{Tr}\left\{  KD_{\ast}\left(  \Gamma_{1}\left(  \mathbf{\varphi
}\right)  \right)  \right\}  . \label{KS4}%
\end{equation}
One should note that

\begin{itemize}
\item for fixed values of $\gamma$ the functionals $\mathcal{F}_{XC}\left(
\Gamma_{1}\left(  \mathbf{\varphi}\right)  \right)  $ and $\mathcal{F}%
_{K}\left(  \Gamma_{1}\left(  \mathbf{\varphi}\right)  \right)  $ in eqs.
$\left(  \ref{KS3}\right)  $ and $\left(  \ref{KS4}\right)  $ are
\emph{constant} for all $\mathbf{\varphi}$ that satisfy $\gamma=\Gamma
_{1}\left(  \mathbf{\varphi}\right)  $, while the functional $\mathcal{F}%
_{T}\left(  \mathbf{\varphi}\right)  $ can have \emph{different} values for
values of $\mathbf{\varphi}$ that correspond to a fixed values of $\gamma$

\item the composition in eq. $\left(  \ref{FO}\right)  $ defines a class of
orbital functionals that have a particular form. One could define orbital
functionals without reference to densities (as done in section ?) and their
form and properties would be quite different.
\end{itemize}

The minimization problem that determines the ground state energy and
space-spin density of a system characterized by a local potential $\mathrm{v}$
is
\begin{equation}
E\left(  \Gamma_{1}\left(  \mathbf{\varphi}_{GS}\right)  \mathbf{,}%
\mathrm{v}\right)  =\min_{\mathbf{\varphi\in}\mathfrak{S}\left(
\mathbf{N}\left(  N\right)  \right)  }\mathcal{E}_{O}\left(  \mathbf{\varphi
,}\mathrm{v}\right)  , \label{KSmin}%
\end{equation}
where
\begin{equation}
\mathcal{E}_{O}\left(  \mathbf{\varphi,}\mathrm{v}\right)  =\mathcal{F}%
_{O}\left(  \mathbf{\varphi}\right)  +f_{\Gamma_{1}\left(  \mathbf{\varphi
}\right)  }\left(  \mathrm{v}\right)  .
\end{equation}
The feasible set $\mathfrak{S}\left(  \mathbf{N}\left(  N\right)  \right)  $
is generated by the constraints
\begin{equation}
\left\langle \left.  \varphi_{j}\right\vert \varphi_{k}\right\rangle
=N_{j}\delta_{jk},\,1\leq j,k\leq r_{1},
\end{equation}
with control parameters $\left\{  \left.  N_{j}\right\vert 1\leq j,\leq
r_{1}\right\}  $ constrained by
\begin{equation}
0\leq N_{j}\leq1,\,1\leq j\leq r_{1};\sum_{\,1\leq j\leq r_{1}}N_{j}=N.
\end{equation}
The GSO minimization problem eq. $\left(  \ref{KSmin}\right)  $ can be
decomposed into two sequential optimizations, first one varies the GSO's
keeping their normalization fixed and determines the saddle point
$\mathcal{L}_{\mathbf{\varphi}}\left(  \mathbf{\varphi}\left(  \mathbf{N}%
\right)  ,\mathbf{\varepsilon}\left(  \mathbf{N}\right)  ,\mathrm{v}\right)  $
of the Lagrangian
\begin{equation}
\mathcal{L}_{\mathbf{\varphi}}\left(  \mathbf{\varphi},\mathbf{\varepsilon
,N,}\mathrm{v}\right)  =\mathcal{E}_{O}\left(  \mathbf{\varphi,}%
\mathrm{v}\right)  -\sum_{1\leq j,k\leq r_{1}}\varepsilon_{jk}\left\{
\overset{}{\left\langle \left.  \varphi_{j}\right\vert \varphi_{k}%
\right\rangle -N_{j}\delta_{jk}}\right\}  ,
\end{equation}
i.e.
\begin{equation}
\mathcal{L}_{\mathbf{\varphi}}\left(  \mathbf{\varphi}\left(  \mathbf{N}%
\right)  ,\mathbf{\varepsilon}\left(  \mathbf{N}\right)  ,\mathrm{v}\right)
=\max_{\mathbf{\varepsilon}}\min_{\mathbf{\varphi}}\mathcal{L}%
_{\mathbf{\varphi}}\left(  \mathbf{\varphi},\mathbf{\varepsilon,N,}%
\mathrm{v}\right)  , \label{KSM1}%
\end{equation}
which leads to optimal values $\mathcal{E}_{O}\left(  \mathbf{\varphi}\left(
\mathbf{N}\right)  \mathbf{,}\mathrm{v}\right)  =$ $\mathcal{L}%
_{\mathbf{\varphi}}\left(  \mathbf{\varphi}\left(  \mathbf{N}\right)
,\mathbf{\varepsilon}\left(  \mathbf{N}\right)  ,\mathbf{N,}\mathrm{v}\right)
$ of the GSO energy functional for fixed values of the control parameters
$\mathbf{N}$. In the second optimization one finds minima of $\mathcal{E}%
_{O}\left(  \mathbf{\varphi}\left(  \mathbf{N}\right)  \mathbf{,}%
\mathrm{v}\right)  $ with respect to a constrained variation of the
normalization control parameters. The constraints
\begin{equation}
0\leq N_{j}\leq1,1\leq j\leq r_{1}%
\end{equation}
on the control parameter $\mathbf{N}$ can be held by setting
\begin{equation}
N_{j}\left(  \mathbf{\theta}\right)  =\sin^{2}\theta_{j},
\end{equation}
and the introduction of a Lagrange parameter $\eta$ in order to enforce the
constraint
\begin{equation}
\sum_{1\leq j\leq r_{1}}N_{j}\left(  \mathbf{\theta}\right)  =N.
\end{equation}
This leads to a Lagrangian
\begin{equation}
\mathcal{L}_{\mathbf{\theta}}\left(  \mathbf{\varphi}\left(  \mathbf{\theta
}\right)  ,\eta,N,\mathrm{v}\right)  =\mathcal{E}_{O}\left(  \mathbf{\varphi
}\left(  \mathbf{\theta}\right)  \mathbf{,}\mathrm{v}\right)  -\eta\left\{
\sum_{\,1\leq j\leq r_{1}}N_{j}\left(  \mathbf{\theta}\right)  -N\right\}  .
\end{equation}
The saddle point $\mathcal{L}_{\mathbf{\theta}}\left(  \mathbf{\varphi}\left(
\mathbf{\theta}_{\#}\left(  N\right)  \right)  ,\eta\left(  N\right)
,\mathrm{v}\right)  $ is given by
\begin{equation}
\mathcal{L}_{\mathbf{\theta}}\left(  \mathbf{\varphi}\left(  \mathbf{\theta
}_{\#}\left(  N\right)  \right)  ,\eta,N,\mathrm{v}\right)  =\max_{\eta}%
\min_{\mathbf{\theta}}\mathcal{L}_{\mathbf{\theta}}\left(  \mathbf{\varphi
}\left(  \mathbf{\theta}\right)  ,\eta,N,\mathrm{v}\right)  ,
\end{equation}
which determines the ground state energy $E\left(  \gamma_{GS}\mathbf{,}%
\mathrm{v}\right)  =\mathcal{E}_{O}\left(  \mathbf{\varphi}\left(
\mathbf{\theta}_{\#}\left(  N\right)  \right)  ,\mathrm{v}\right)  .$The
ground state density can then expanded as
\begin{equation}
\gamma_{GS}\left(  \mathbf{x}\right)  =\sum_{1\leq j\leq r_{1}}\left\vert
\overset{}{\left[  \varphi_{j}\left(  \mathbf{\theta}_{\#}\left(  N\right)
\right)  \right]  \left(  \mathbf{x}\right)  }\right\vert ^{2}.
\end{equation}


The saddle points of $\mathcal{L}_{\mathbf{\varphi}}\left(  \mathbf{\varphi
},\mathbf{\varepsilon,N,}\mathrm{v}\right)  $ are characterized by the
necessary conditions
\begin{subequations}
\begin{align}
\frac{\delta\mathcal{L}_{\mathbf{\varphi}}\left(  \mathbf{\varphi
},\mathbf{\varepsilon,N,}\mathrm{v}\right)  }{\delta\varphi_{j}^{\ast}}  &
=0;\,1\leq j\leq r_{1},\label{KSC1}\\
\frac{\partial\mathcal{L}_{\mathbf{\varphi}}\left(  \mathbf{\varphi
},\mathbf{\varepsilon,N,}\mathrm{v}\right)  }{\partial\varepsilon_{jk}}  &
=0;\,1\leq j,k\leq r_{1} \label{KSC2}%
\end{align}
(Appendix \ref{ComplexConjugate} describes why it is sufficient to consider
derivatives wrt $\varphi_{j}^{\ast}).$ The condition eq. $\left(
\ref{KSC1}\right)  $ implies that at an optimal point the matrix of Lagrange
multipliers $\mathbf{\varepsilon}\left(  \mathbf{N}\right)  $ is hermitian
i.e.
\end{subequations}
\begin{equation}
\varepsilon\left(  \mathbf{N}\right)  _{jk}=\varepsilon\left(  \mathbf{N}%
\right)  _{kj}^{\ast},\quad1\leq j,k\leq r_{1}%
\end{equation}
and
\begin{equation}
\frac{\delta\mathcal{E}_{O}\left(  \mathbf{\varphi}\left(  \mathbf{N}\right)
,\mathbf{\varepsilon}\left(  \mathbf{N}\right)  ,\mathrm{v}\right)  }%
{\delta\varphi_{j}^{\ast}}-\sum_{1\leq k\leq r_{1}}\varepsilon_{jk}\left(
\mathbf{N}\right)  \left\vert \varphi_{k}\left(  \mathbf{N}\right)
\right\rangle =0;\,1\leq j\leq r_{1}.
\end{equation}
Which after noting that if one factors the energy functional through the
density, (denoted by $\mathcal{E}_{O}\left(  \gamma\left(  \mathbf{\varphi
}\right)  \mathbf{,}\mathrm{v}\right)  $ in contrast to $\mathcal{E}%
_{O}\left(  \mathbf{\varphi,}\mathrm{v}\right)  $), one obtains
\begin{align}
\frac{\delta\mathcal{E}_{O}\left(  \gamma\left(  \mathbf{\varphi}\right)
\mathbf{,}\mathrm{v}\right)  }{\delta\varphi_{j}^{\ast}\left(  \mathbf{x}%
\right)  }  &  =\int\frac{\delta\mathcal{E}_{O}\left(  \gamma\left(
\mathbf{\varphi}\right)  \mathbf{,}\mathrm{v}\right)  }{\delta\gamma\left(
\mathbf{x}^{\prime}\right)  }\frac{\delta\gamma\left(  \mathbf{x}^{\prime
}\right)  }{\delta\varphi_{j}^{\ast}\left(  \mathbf{x}\right)  }%
d\mathbf{x}^{\prime}\nonumber\\
&  =\int\frac{\delta\mathcal{E}_{O}\left(  \gamma\left(  \mathbf{\varphi
}\right)  \mathbf{,}\mathrm{v}\right)  }{\delta\gamma\left(  \mathbf{x}%
\right)  }\varphi_{j}\left(  \mathbf{x}^{\prime}\right)  \delta\left(
\mathbf{x-x}^{\prime}\right)  d\mathbf{x}^{\prime}=\frac{\delta\mathcal{E}%
_{O}\left(  \gamma\left(  \mathbf{\varphi}\right)  \mathbf{,}\mathrm{v}%
\right)  }{\delta\gamma\left(  \mathbf{x}\right)  }\varphi_{j}\left(
\mathbf{x}\right)
\end{align}
that gives%
\begin{equation}
\frac{\delta\mathcal{E}_{O}\left(  \left[  \gamma\left(  \mathbf{\varphi
}\right)  \right]  \mathbf{\left(  \mathbf{N}\right)  ,}\mathrm{v}\right)
}{\delta\gamma}\left\vert \varphi_{j}\left(  \mathbf{N}\right)  \right\rangle
-\sum_{1\leq k\leq r_{1}}\varepsilon_{jk}\left(  \mathbf{N}\right)  \left\vert
\varphi_{k}\left(  \mathbf{N}\right)  \right\rangle =0;\,1\leq j\leq r_{1}.
\end{equation}
The operator $\frac{\delta\mathcal{E}_{O}\left(  \left[  \gamma\left(
\mathbf{\varphi}\right)  \right]  \mathbf{\left(  \mathbf{N}\right)
,}\mathrm{v}\right)  }{\delta\gamma}$ is a multiplication operator in the
coordinate representation given by
\begin{equation}
\frac{\delta\mathcal{E}_{O}\left(  \left[  \gamma\left(  \mathbf{\varphi
}\right)  \right]  \mathbf{\left(  \mathbf{N}\right)  ,}\mathrm{v}\right)
}{\delta\gamma}\left\vert \varphi_{j}\left(  \mathbf{N}\right)  \right\rangle
=\left\vert \frac{\delta\mathcal{E}_{O}\left(  \left[  \gamma\left(
\mathbf{\varphi}\right)  \right]  \mathbf{\left(  \mathbf{N}\right)
,}\mathrm{v}\right)  }{\delta\gamma}\varphi_{j}\left(  \mathbf{N}\right)
\right\rangle ,
\end{equation}
where
\begin{equation}
\left[  \frac{\delta\mathcal{E}_{O}\left(  \left[  \gamma\left(
\mathbf{\varphi}\right)  \right]  \mathbf{\left(  \mathbf{N}\right)
,}\mathrm{v}\right)  }{\delta\gamma}\varphi_{j}\left(  \mathbf{N}\right)
\right]  \left(  \mathbf{x}\right)  =\left(  \frac{\delta\mathcal{E}%
_{O}\left(  \left[  \gamma\left(  \mathbf{\varphi}\right)  \right]
\mathbf{\left(  \mathbf{N}\right)  ,}\mathrm{v}\right)  }{\delta\gamma\left(
\mathbf{x}\right)  }\right)  \left(  \left[  \varphi_{j}\left(  \mathbf{N}%
\right)  \right]  \left(  \mathbf{x}\right)  \right)  .
\end{equation}
We thus obtain the intermediate GSO equations that correspond to condition eq.
$\left(  \ref{KSC1}\right)  $
\begin{equation}
\frac{\delta\mathcal{E}_{O}\left(  \left[  \gamma\left(  \mathbf{\varphi
}\right)  \right]  \mathbf{\left(  \mathbf{N}\right)  ,}\mathrm{v}\right)
}{\delta\gamma}\left\vert \varphi_{j}\left(  \mathbf{N}\right)  \right\rangle
=\sum_{1\leq k\leq r_{1}}\varepsilon_{jk}\left(  \mathbf{N}\right)  \left\vert
\varphi_{k}\left(  \mathbf{N}\right)  \right\rangle ,\,\varepsilon_{jk}\left(
\mathbf{N}\right)  =\varepsilon_{kj}^{\ast}\left(  \mathbf{N}\right)  .
\label{FockEq}%
\end{equation}
One should note that the norm of the vectors $\left\{  \left\vert \varphi
_{j}\left(  \mathbf{N}\right)  \right\rangle \right\}  $ are the occupation
numbers and eq. $\left(  \ref{FockEq}\right)  $ is not an eigenvalue equation
unless $\mathcal{E}_{O}\left(  \mathbf{\varphi}\left(  \mathbf{N}\right)
,\mathrm{v}\right)  $ is invariant to unitary mixing of $\mathbf{\varphi
}\left(  \mathbf{N}\right)  $.

From sensitivity theory the rate of change of the orbital functional energy
with respect to changes in occupation numbers, which in this case are the
control parameters $\mathbf{N}$, is
\begin{equation}
\frac{\partial\mathcal{E}_{O}\left(  \varphi\mathbf{\left(  \mathbf{N}\right)
},\mathrm{v}\right)  }{\partial N_{j}}=\varepsilon_{jj}\left(  \mathbf{N}%
\right)  , \label{Janak}%
\end{equation}
which is Janak's Theorem extended to non normalized GSO's and \emph{has been
obtained} without consideration of the expansion of the density in terms of
GSO's. If we take that expansion into account we obtain this relationship as
the expectation value of the operator $\frac{\delta\mathcal{E}_{O}\left(
\left[  \gamma\left(  \mathbf{\varphi}\right)  \right]  \mathbf{\left(
\mathbf{N}\right)  ,}\mathrm{v}\right)  }{\delta\gamma}$ i.e.
\begin{equation}
\frac{\partial\mathcal{E}_{O}\left(  \left[  \gamma\left(  \mathbf{\varphi
}\right)  \right]  \mathbf{\left(  \mathbf{N}\right)  },\mathrm{v}\right)
}{\partial N_{j}}=\left\langle \tilde{\varphi}_{j}\left(  \mathbf{N}\right)
\left\vert \frac{\delta\mathcal{E}_{O}\left(  \left[  \gamma\left(
\mathbf{\varphi}\right)  \right]  \mathbf{\left(  \mathbf{N}\right)
,}\mathrm{v}\right)  }{\delta\gamma}\right.  \tilde{\varphi}_{j}\left(
\mathbf{N}\right)  \right\rangle =\tilde{\varepsilon}_{jj}\left(
\mathbf{N}\right)  =\varepsilon_{jj}\left(  \mathbf{N}\right)  ,
\end{equation}
(see Appendix \ref{GSO&Occ}).

After noting that
\begin{equation}
\frac{\partial\mathcal{L}_{\mathbf{\theta}}\left(  \mathbf{\varphi}\left(
\mathbf{\theta}\right)  ,\eta,N,\mathrm{v}\right)  }{\partial\theta_{j}}%
=\frac{\partial\mathcal{E}_{O}\left(  \varphi\left(  \mathbf{\theta}\right)
,\mathrm{v}\right)  }{\partial\theta_{j}}-\eta\sin\left\{  2\theta
_{j}\right\}  =\left\{  \frac{\partial\mathcal{E}_{O}\left(  \varphi\left(
\mathbf{N}\left(  \mathbf{\theta}\right)  \right)  ,\mathrm{v}\right)
}{\partial N_{j}}-\eta\right\}  \sin\left\{  2\theta_{j}\right\}  \,;\,1\leq
j\leq r_{1},
\end{equation}
one observes that the stationarity conditions eq. $\left(  \ref{KSC2}\right)
$
\begin{equation}
\frac{\partial\mathcal{L}_{\mathbf{\theta}}\left(  \mathbf{\varphi}\left(
\mathbf{\theta}\left(  N\right)  \right)  ,\eta\mathbf{,}N\mathbf{,}%
\mathrm{v}\right)  }{\partial\theta_{j}}=0\text{ }%
\end{equation}
for the outer optimization give%
\begin{equation}
\frac{\partial\mathcal{E}_{O}\left(  \varphi\left(  \mathbf{N}\left(
\mathbf{\theta}\left(  N\right)  \right)  \right)  ,\mathrm{v}\right)
}{\partial N_{j}}=\eta\left(  N\right)  =\varepsilon_{jj}\left(
\mathbf{\mathbf{N}}\left(  \mathbf{\theta}\left(  N\right)  \right)  \right)
\label{out1}%
\end{equation}
or%
\begin{equation}
\sin\left\{  2\theta_{j}\left(  N\right)  \right\}  =0. \label{out2}%
\end{equation}
The stationary condition eq. $\left(  \ref{out2}\right)  $ gives
\begin{equation}
\sin\left\{  2\theta_{j}\left(  N\right)  \right\}  =0\Rightarrow\theta
_{j}\left(  N\right)  =0\text{ or }\frac{\pi}{2}\Rightarrow N_{j}\left(
\mathbf{\theta}\left(  N\right)  \right)  =0\text{ or }1
\end{equation}
and corresponds to values on the boundary of the feasable set occupation
numbers. Hence
\begin{equation}
\frac{\partial\mathcal{E}_{O}\left(  \mathbf{\varphi}\left(  \mathbf{N}\left(
\mathbf{\theta}_{\#}\left(  N\right)  \right)  \right)  \mathbf{,}%
\mathrm{v}\right)  }{\partial N_{j}}=\eta\left(  N\right)  \text{ or }%
N_{j}\left(  \mathbf{\theta}_{\#}\left(  N\right)  \right)  =0\text{ or
}1;\,1\leq j\leq r_{1}.
\end{equation}


The preceeding relationships are true for \emph{any} orbital theory, where the
occupation numbers are constariend to lie between zero and one and add up to
the number of electrons in the system. By considering sensitivity theory one
also obtains
\begin{equation}
\frac{\partial\mathcal{E}_{O}\left(  \mathbf{\varphi}\left(  \mathbf{\theta
}\left(  N\right)  \right)  \mathbf{,}\mathrm{v}\right)  }{\partial N}%
=\eta\left(  N\right)
\end{equation}
and $\eta\left(  N\right)  $ can be interpreted as the chemical potential of
the interacting molecular system.

\subsubsection{The Kohn-Sham Density Functional}

The minima of $\mathcal{E}_{O}\left(  \mathbf{\varphi,}\mathrm{v}\right)  ,$
for both the inner and outer minimizations, are highly degenerate e.g. there
are many minimizers $\mathbf{\varphi}\left(  \mathbf{\theta}_{\#}\left(
N\right)  \right)  $ that correspond to the constrained minimum value of
$\mathcal{E}_{O}\left(  \mathbf{\varphi}\left(  \mathbf{\theta}_{\#}\left(
N\right)  \right)  \mathbf{,}\mathrm{v}\right)  ,$ even if the ground state
density $\gamma_{GS}$ is unique$,$ as
\begin{equation}
\mathcal{E}_{O}\left(  \mathbf{\varphi,}\mathrm{v}\right)  =const.\,\forall
\mathbf{\varphi\in}\mathfrak{S}\left(  \mathbf{N}\left(  N\right)
\mathbf{,\gamma}\right)  ,
\end{equation}
where $\mathfrak{S}\left(  \mathbf{N}\left(  N\right)  \mathbf{,\gamma
}\right)  $ is the equivalence class
\begin{equation}
\mathfrak{S}\left(  \mathbf{N}\left(  N\right)  \mathbf{,\gamma}\right)
=\mathfrak{S}\left(  \mathbf{N}\left(  N\right)  \right)  \cap\mathfrak{S}%
\left(  \mathbf{N,\gamma}\right)
\end{equation}
of sets of orthogonal GSO's that produce densities that describe systems that
contain on average $N$-electrons. One can remove this degeneracy by
constraining $\mathbf{\varphi}$ to be a specific representative of the
equivalence class $\mathfrak{S}\left(  \mathbf{N}\left(  N\right)
\mathbf{,\gamma}\right)  .$ In order to obtain such a representative one
considers the GSO's formed by the NGSO's and occupation numbers of the FORDO
of the $N$-electron ground state $D_{0}\left(  \gamma\right)  =A_{0}\left(
\gamma\right)  ^{2}$ of a model independent particle $N$-electron system$,$
that has a fixed density $\gamma.$ This corresponds to finding the minimizers
of the universal kinetic energy functional $K_{U}$
\begin{equation}
K_{U}\left(  A\right)  =\operatorname{Tr}\left\{  KA^{2}\right\}  ,
\end{equation}
subject to the constraint
\begin{equation}
\Xi_{2}\left(  A\right)  =\left(  A\left\vert \widehat{\Phi^{\dagger}\Phi
}A\right.  \right)  =\gamma,\,A\in\mathcal{B}_{2sa}\left(  \mathcal{H}%
^{N}\right)  .
\end{equation}
The constrained minima $K_{U}\left(  D_{0}\left(  \gamma\right)  \right)
=\operatorname{Tr}\left\{  KD_{0}\left(  \gamma\right)  \right\}
=\operatorname{Tr}\left\{  KA_{0}\left(  \gamma\right)  ^{2}\right\}  $ is
given as an unconstrained saddle point $\left(  A_{0}\left(  \gamma\right)
,\lambda_{0}\left(  \gamma\right)  \right)  ,$
\begin{equation}
\mathcal{L}_{K}\left(  A_{0}\left(  \gamma\right)  ,\lambda_{0}\left(
\gamma\right)  ,\gamma\right)  =\max_{\lambda}\min_{A}\mathcal{L}_{K}\left(
A,\lambda,\gamma\right)
\end{equation}
of the Lagrangian
\begin{equation}
\mathcal{L}_{K}\left(  A,\lambda,\gamma\right)  =K_{U}\left(  A\right)
-\Omega_{A}\left(  \lambda\right)  +\left\langle \left.  \lambda\right\vert
\gamma\right\rangle =\left(  A\left\vert \left\{  \widehat{\left\{
K-\Omega\left(  \lambda\right)  \right\}  }+\left\langle \left.
\lambda\right\vert \gamma\right\rangle \hat{I}\right\}  \right.  A\right)  .
\end{equation}
If a bound state minimum exists for some local one body potential $\lambda
_{0}\left(  \gamma\right)  $ one has
\begin{equation}
\widehat{\left\{  K-\Omega\left(  \lambda_{0}\left(  \gamma\right)  \right)
\right\}  }\left\vert A_{0}\left(  \gamma\right)  \right)  =0,
\end{equation}
where $A_{0}\left(  \gamma\right)  \in\mathcal{B}_{2sa}\left(  \mathcal{H}%
^{N}\right)  $. This gives a stationary state $D_{0}\left(  \gamma\right)
=A_{0}\left(  \gamma\right)  ^{2}$ of the independent particle system with
Hamiltonian $\left\{  K-\Omega\left(  \lambda_{0}\left(  \gamma\right)
\right)  \right\}  $ i.e.
\begin{equation}
\left[  \left\{  K-\Omega\left(  \lambda_{0}\left(  \gamma\right)  \right)
\right\}  ,D_{0}\left(  \gamma\right)  \right]  =0,
\end{equation}
which in turn defines a FORDO $D_{0}^{1}\left(  \gamma\right)  $ given by (see
Appendix \ref{Denmat})
\begin{equation}
D_{0}^{1}\left(  \gamma\right)  =NL_{N}^{1}\left(  D_{0}\left(  \gamma\right)
\right)  .
\end{equation}
As $K-\Omega\left(  \lambda_{0}\left(  \gamma\right)  \right)  $ is an one
electron operator the state $D_{0}\left(  \gamma\right)  $ must be a mixture
of IPS's or a pure IPS and thus determined by the eigenvalues $\left\{
N_{0j}\left(  \gamma\right)  \right\}  $ and eigenvectors $\left\{  \left\vert
\widetilde{\varphi_{0j}\left(  \gamma\right)  }\right\rangle \right\}  $ of
$D_{0}^{1}\left(  \gamma\right)  $. Expanding $D_{0}^{1}\left(  \gamma\right)
$ in terms of normalized NGSO's $\left\{  \left.  \left\vert \widetilde
{\varphi_{0j}}\right\rangle \right\vert 1\leq j\leq r_{1}\right\}  $ gives
\begin{equation}
D_{0}^{1}\left(  \gamma\right)  =\sum_{1\leq j\leq r_{1}}N_{0j}\left(
\gamma\right)  \left\vert \widetilde{\varphi_{0j}\left(  \gamma\right)
}\right\rangle \left\langle \widetilde{\varphi_{0j}\left(  \gamma\right)
}\right\vert ,
\end{equation}
(where in contrast to eq. $\left(  \ref{Gamma}\right)  $ $r_{1}\geq N$ ),
which after renormalizing the NGSO's to the occupation numbers i.e.
\begin{equation}
\left\vert \varphi_{0j}\left(  \gamma\right)  \right\rangle =N_{0j}\left(
\gamma\right)  ^{\frac{1}{2}}\left\vert \widetilde{\varphi_{0j}\left(
\gamma\right)  }\right\rangle ,1\leq j\leq r_{1}%
\end{equation}
results in
\begin{equation}
\gamma\left(  \mathbf{x}\right)  =\sum_{1\leq j\leq r_{1}}\left\vert \left[
\varphi_{0j}\left(  \gamma\right)  \right]  \left(  \mathbf{x}\right)
\right\vert ^{2}. \label{gammaKS}%
\end{equation}
If the constrained minimum of $K_{U}$ is non degenerate then there is a unique
correspondence, up to unitary mixing of degenerate eigenvectors of $D_{0}%
^{1}\left(  \gamma\right)  ,$ between $\gamma$ and $\left\{  \mathbf{\varphi
}_{0}\right\}  ,$ thus modulo this unitary mixing one can define an inverse
functional $\Gamma_{0}^{-1}$ as
\begin{equation}
\mathbf{\varphi}_{0}=\Gamma_{0}^{-1}\left(  \gamma\right)  .
\end{equation}
If the constrained minimum of $K_{U}$ is degenerate, e.g. corresponding to
FORDO's $\left\{  \left.  D_{0\kappa}^{1}\left(  \gamma\right)  \right\vert
1\leq\kappa\leq m\right\}  $ of the stationary states $\left\{  \left.
D_{0\kappa}\left(  \gamma\right)  \right\vert 1\leq\kappa\leq m\right\}  $ of
$K-\Omega\left(  \lambda_{0}\left(  \gamma\right)  \right)  $ then as
\begin{equation}
\left[  \left\{  K-\Omega\left(  \lambda_{0}\left(  \gamma\right)  \right)
\right\}  ,D_{0\kappa}\left(  \gamma\right)  \right]  =\left[  \left\{
K-\Omega\left(  \lambda_{0}\left(  \gamma\right)  \right)  \right\}
,D_{0\kappa}^{1}\left(  \gamma\right)  \right]  =0,
\end{equation}
the space of degenerate states corresponds to the same set of GSO's $\left\{
\left\vert \mathbf{\varphi}_{0}\left(  \gamma\right)  \right\rangle \right\}
$ but normalized to different occupation numbers $\left\{  \mathbf{N}%
_{0}\left(  \gamma\right)  \right\}  ,$ where $\mathbf{N}_{0}\left(
\gamma\right)  $ are convex combinations of $\left\{  \left.  \mathbf{N}%
_{0\kappa}\left(  \gamma\right)  \right\vert 1\leq\kappa\leq m\right\}  $.
Hence if degeneracies occur these are reflected in multiple optimal values for
$\mathbf{\mathbf{N}}_{0}$ for fixed $\left\{  \mathbf{\varphi}_{0}\left(
\gamma\right)  \right\}  \mathbf{.}$

The minimal values of the universal kinetic functional $K_{U}\left(
D_{0}\left(  \gamma\right)  \right)  $ produce a kinetic energy functional
$\mathcal{F}_{T}$ defined on the representatives $\left\{  \mathbf{\varphi
}_{0}\right\}  $ of the equivalence classes $\mathfrak{S}\left(
\mathbf{N}\left(  N\right)  \mathbf{,}\gamma\right)  $ whose values are given
by
\begin{equation}
\mathcal{F}_{T}\left(  \mathbf{\varphi}_{0}\right)  =K_{U}\left(  D_{0}\left(
\Gamma_{1}\left(  \mathbf{\varphi}_{0}\right)  \right)  \right)  .
\label{KEFunc}%
\end{equation}
The value of $\mathcal{F}_{T}$ is constant on the set of GSO's formed from the
unitary mixing and the convex combinations of degenerate states just
described. This allows one to define the Kohn-Sham functional $\mathcal{F}%
_{KS}\left(  \mathbf{\varphi}_{0}\right)  $ as a specific decomposition of the
HK energy functional $\mathcal{F}_{HK}\left(  \gamma\right)  $ in terms of
GSO's as
\begin{equation}
\mathcal{F}_{KS}\left(  \mathbf{\varphi}_{0}\right)  =\mathcal{F}_{T}\left(
\mathbf{\varphi}_{0}\right)  +\mathcal{F}_{C}\left(  \mathbf{\varphi}%
_{0}\right)  +\mathcal{F}_{KSXC}\left(  \mathbf{\varphi}_{0}\right)  ,
\end{equation}
where
\begin{equation}
\mathcal{F}_{KS}\left(  \mathbf{\varphi}_{0}\right)  =\mathcal{F}_{HK}\left(
\Gamma\left(  \mathbf{\varphi}_{0}\right)  \right)  .
\end{equation}


The relationships described in subsection \ref{GSO} derived for stationary
points $\mathbf{\varphi}\left(  \mathbf{N}\left(  \mathbf{\theta}\right)
\right)  ,$ $\mathbf{\varphi}\left(  \mathbf{N}\left(  \mathbf{\theta}%
_{\#}\left(  N\right)  \right)  \right)  $ are valid for any points contained
in the equivalence classes that these points belong to, in particular
$\mathbf{\varphi}_{0}\left(  \mathbf{N}\left(  \mathbf{\theta}\right)
\right)  $ and $\mathbf{\varphi}_{0}\left(  \mathbf{N}\left(  \mathbf{\theta
}_{\#}\left(  N\right)  \right)  \right)  .$

The reference $N$-electron IPS, $D_{0}\left(  \mathbf{N}\left(  \mathbf{\theta
}_{\#}\left(  N\right)  \right)  \right)  $ formed from the GSO's
$\mathbf{\varphi}_{0}$ has two energies associated with it. One, $E_{0}\left(
\mathbf{\varphi}_{0}\left(  \mathbf{N}\left(  \mathbf{\theta}_{\#}\left(
N\right)  \right)  \right)  \right)  $, is the ground state energy of the
Hamiltonian%
\begin{equation}
H_{0}=K+\mathrm{v},
\end{equation}
which can be expressed as
\begin{equation}
E_{0}\left(  \mathbf{\varphi}_{0}\left(  \mathbf{N}\left(  \mathbf{\theta
}_{\#}\left(  N\right)  \right)  \right)  \right)  =\sum_{1\leq j\leq r_{1}%
}\left\langle \tilde{\varphi}_{0j}\left\vert \left(  K+\mathrm{v}\right)
\right.  \tilde{\varphi}_{0j}\right\rangle N_{j}\left(  \mathbf{\theta}%
_{\#}\left(  N\right)  \right)  .
\end{equation}
The other is the energy of the Hamiltonian
\begin{equation}
H=H_{0}+V_{2}%
\end{equation}
wrt the IPS $D_{0}\left(  \mathbf{N}\left(  \mathbf{\theta}_{\#}\left(
N\right)  \right)  \right)  $ given by
\begin{equation}
E_{IP}\left(  D_{0}\left(  \mathbf{N}\left(  \mathbf{\theta}_{\#}\left(
N\right)  \right)  \right)  \right)  =\operatorname{Tr}\left\{  HD_{0}\left(
\mathbf{N}\left(  \mathbf{\theta}_{\#}\left(  N\right)  \right)  \right)
\right\}  .
\end{equation}
It should be noted that this value of the energy functional $E_{IP}$ is
\emph{not }optimized in any way and is \emph{not }related in any simple
fashion to the ground state energy $E\left(  \gamma_{GS}\mathbf{,}%
\mathrm{v}\right)  .$

The classical Janaks Theorem is a special case of eq. $\left(  \ref{Janak}%
\right)  $
\begin{equation}
\frac{\partial\mathcal{E}_{O}\left(  \mathbf{\varphi}_{0}\left(
\mathbf{N}\left(  \mathbf{\theta}\left(  N\right)  \right)  \right)
,\mathrm{v}\right)  }{\partial N_{j}}\equiv\frac{\partial\mathcal{E}%
_{KS}\left(  \mathbf{\varphi}_{0}\left(  \mathbf{N}\left(  \mathbf{\theta
}\left(  N\right)  \right)  \right)  ,\mathrm{v}\right)  }{\partial N_{j}%
}=\varepsilon_{jj}\left(  \mathbf{N}\left(  \mathbf{\theta}\left(  N\right)
\right)  \right)  , \label{ClassJanak}%
\end{equation}
where we have constrained the GSO's and occupation numbers to be the ones that
minimize the universal kinetic energy, $K_{U},$ thus producing a model IPS.
One must, however, note that the Lagrange parameters $\left\{  \varepsilon
_{jj}\left(  \mathbf{N}\left(  \mathbf{\theta}\left(  N\right)  \right)
\right)  \right\}  $ in this equation \emph{are not }associated with the
minimization of the universal kinetic energy.\emph{ }

The KS energy functional is defined in terms of GSO's and occupation numbers
and is implicitly defined by the preceding constrained optimization
procedures, that factor the functional through the density. If one did not
consider this factorization and directly optimized the GSO's and occupation
numbers, i.e. use a FORDO functional one would determine the same ground
state. In the next section we compare and contrast the relationships obtained
for these functionals with ones obtained for the explicitly defined AGP
functional. These functionals are defined on the space of correlated states
formed by the non linear manifold of AGP\ states, where the energy has been
expressed as an explicit functional of GSO's and parameters that determine
occupation numbers \cite{WeinerOrtiz05}. The minimization of this explicit
functional can then be compared to the minimization of the implicit
KS\ functional restricted to the AGP\ manifold.

\section{AGP Energy Functionals \label{AGP}}

We briefly review some of the results found in \cite{weinerortiz05} describing
AGP states and their associated energy functionals. In \cite{weinerortiz05}
orbital and occupation number optimization equations were derived and it was
displayed how the $N$-electron total energy could be expressed in terms of
ionization potentials, kinetic energies and eigenvalues of a generalized Fock
operator associated with GSO's. In contrast to the KS approach, the occupation
numbers were not constrained to refer to IPS's, but allowed to vary over the
much larger domain of states with doubly degenerate state occupation numbers.
It has been argued while this set does not contain all possible states it
contains a large percentage of states important for molecular systems
\cite{Kramers Deg}. The density matrix energy functional in DMFT allows
unrestricted variation of occupation numbers (and thus all states), but unlike
the AGP functional (but in common with the HK\ and KS density functionals) it
is invariant to the phases of the GSO's i.e. the group $U\left(  1\right)  .$

\subsection{Canonical expansion of an AGP state}

\bigskip The $2M$-electron AGP state, $\left\vert g^{M}\right\rangle $, is the
$M^{th}$ antisymmetric power of the $2$-electron state described by the
geminal
\begin{equation}
\left\vert g\right\rangle =\sum_{1\leq j\leq s}n_{j}^{\frac{1}{2}}\left\vert
\psi_{j}\psi_{j+s}\right\rangle . \label{geminal}%
\end{equation}
given by
\begin{equation}
\left\vert g^{M}\right\rangle =\left(  S_{M}\right)  ^{-\frac{1}{2}}%
\sum_{1\leq j_{1}<\cdots<j_{M}\leq s}n_{j_{1}}^{\frac{1}{2}}\cdots n_{j_{s}%
}^{\frac{1}{2}}\left\vert \psi_{j_{1}}\psi_{j_{1}+s}\cdots\psi_{j_{M}}%
\psi_{j_{M}+s}\right\rangle . \label{agp}%
\end{equation}
$\left\{  \left.  n_{j}^{\frac{1}{2}}\right\vert 1\leq j\leq s\right\}  $ are
the positive square roots of the occupation numbers $\left\{  \left.
n_{j}\right\vert 1\leq j\leq s\right\}  $ of the FORDO of the two electron
AGP\ state density operator and $\left\{  \left.  \left\vert \psi
_{j}\right\rangle \right\vert 1\leq j\leq2s\right\}  $ are the Canonical
General Spin Orbitals (CGSO's) of $\left\vert g\right\rangle .$ The
normalization factor, $\left(  S_{M}\right)  ^{-\frac{1}{2}}$, is given in
terms of the elementary symmetric function
\begin{equation}
S_{M}=\sum_{1\leq j_{1}<\cdots<j_{M}\leq s}n_{j_{1}}\cdots n_{j_{M}},
\end{equation}
and the geminals $\left\vert g\right\rangle $ represent normalized two
electron states so that
\begin{equation}
\sum_{1\leq j\leq s}n_{j}=1.
\end{equation}
Note that by definition
\begin{equation}
S_{0}=1.
\end{equation}
It is evident from eq.\ $\left(  \ref{geminal}\right)  $ that the geminal and
geminal power state are at the least invariant to the group
\begin{equation}
SU\left(  2\right)  ^{s}=\overset{s\text{-factors}}{\overbrace{\,SU\left(
2\right)  \times\cdots\times SU\left(  2\right)  }},
\end{equation}
that is formed by $2\times2$ \emph{special} unitary transformations of the
paired CGSOs $\left\{  \left.  \left\vert \psi_{j}\right\rangle ,\left\vert
\psi_{j+s}\right\rangle \right\vert 1\leq j\leq s\right\}  ,$ while if
degeneracies exist amongst the $\left\{  \left.  n_{j}\right\vert 1\leq j\leq
s\right\}  $ this invariance group is larger \cite{AGPCoherent}. The canonical
CSGOs are unique up to transformations belonging this invariance group.

\subsection{AGP Energy Functionals}

The energy of the AGP state can be expressed as a constrained functional
$\mathcal{E}_{1}$ of geminal occupation numbers and orthonormal CGSO's, by
\emph{ }
\begin{align}
\mathcal{E}_{1}\left(  \mathbf{\tilde{\varphi},n}\right)   &  =\frac
{\mathcal{H}\left(  \mathbf{\tilde{\varphi},n}\right)  }{S_{M}\left(
\mathbf{n}\right)  }\nonumber\\
&  =\frac{1}{S_{M}\left(  \mathbf{n}\right)  }\left\{  \sum_{1\leq j\leq
s}n_{j}S_{M-1}\left(  \jmath\right)  \left\{  \left\langle \tilde{\varphi}%
_{j}\left\vert h_{1}\right.  \tilde{\varphi}_{j}\right\rangle +\left\langle
\tilde{\varphi}_{j+s}\left\vert h\right.  \tilde{\varphi}_{j+s}\right\rangle
+\left\langle \tilde{\varphi}_{j}\tilde{\varphi}_{j+s}||\tilde{\varphi}%
_{j}\tilde{\varphi}_{j+s}\right\rangle \right\}  \right. \nonumber\\
&  +\sum_{1\leq j,k\leq s}n_{j}^{\frac{1}{2}}n_{k}^{\frac{1}{2}}S_{M-1}\left(
\jmath k\right)  \left\langle \tilde{\varphi}_{j}\tilde{\varphi}_{j+s}%
||\tilde{\varphi}_{k}\tilde{\varphi}_{k+s}\right\rangle \left.  +\sum_{1\leq
j<k\leq s}n_{i}n_{j}S_{M-2}\left(  \jmath k\right)  \left\langle
\tilde{\varphi}_{j}\tilde{\varphi}_{k}|||\tilde{\varphi}_{j}\tilde{\varphi
}_{k}\right\rangle \right\}  \label{energy}%
\end{align}
or as an unconstrained functional $\mathcal{E}_{2}$ of one electron general
linear group parameters $\left\{  \mathbf{\eta},\mathbf{\lambda}\right\}  $
as
\begin{equation}
\mathcal{E}_{2}\left(  \mathbf{\eta},\mathbf{\lambda}\right)  =\frac
{\mathcal{H}\left(  \mathbf{\eta},\mathbf{\lambda}\right)  }{S_{M}\left(
\mathbf{\eta}\right)  }=\frac{\left\langle \left.  g\left(  \mathbf{\eta
},\mathbf{\lambda}\right)  ^{M}\right\vert Hg\left(  \mathbf{\eta
},\mathbf{\lambda}\right)  ^{M}\right\rangle }{S_{M}\left(  \mathbf{\eta
}\right)  }, \label{gp_funct}%
\end{equation}
where%
\begin{equation}
S_{M}\left(  \mathbf{\eta}\right)  =\left\langle \left.  k\left(
\mathbf{\eta}\right)  g_{0}^{M}\right\vert k\left(  \mathbf{\eta}\right)
g_{0}^{M}\right\rangle =\sum_{1\leq j_{1}<\cdots<j_{M}\leq s}e^{2\left(
\eta_{j_{1}}+\cdots+\eta_{j_{M}}\right)  }n_{0j_{1}}\cdots n_{0j_{N}}%
\end{equation}
and $\mathbf{\lambda\in}\mathbb{R}^{r_{1}\left(  r_{1}-1\right)  /2}.$

\subsection{AGP Lagrangians\label{Lagran}}

In order to minimize the energy of the $2M$-electron state the over the
manifold of AGP\ states using the functional $\mathcal{E}_{1}$ one needs to
either constrain the GSO's $\mathbf{\varphi}$ to be orthonormal, in which case
we denote the energy functional as $\mathcal{E}_{1}\left(  \mathbf{\tilde
{\varphi},n}\right)  ,$ or orthogonal and normalized to the two electron
geminal occupation numbers, in which case we denote the energy functional as
$\mathcal{E}_{1}\left(  \mathbf{\varphi}\right)  ,$ that satisfy the
constraints
\begin{align}
0  &  \leq n_{j}\leq1;1\leq j\leq s\nonumber\\
&  \sum_{1\leq j\leq s}n_{j}=1,
\end{align}
i.e. $\mathbf{n\in}\mathfrak{N}\left(  1\right)  \subset\mathbb{R}^{s}.$ This
corresponds to treating the geminal occupation numbers either as independent
variables or control parameters in an analogous fashion to section \ref{GSO}.
To this end one can define Lagrangians
\begin{equation}
\mathcal{L}_{I}\left(  \mathbf{\tilde{\varphi},\theta},\mathbf{\tilde
{\varepsilon},}\eta\right)  =\mathcal{E}_{1I}\left(  \mathbf{\tilde{\varphi
},\mathbf{\theta}}\right)  -\sum_{1\leq j,k\leq2s}\tilde{\varepsilon}%
_{kj}\left(  \left\langle \left.  \tilde{\varphi}_{j}\right\vert
\tilde{\varphi}_{k}\right\rangle -\delta_{jk}\right)  -\eta\left\{
\sum_{1\leq j\leq2s}n_{j}\mathbf{\left(  \mathbf{\theta}\right)  }-1\right\}
\label{LI1}%
\end{equation}
or%
\begin{equation}
\mathcal{L}_{C}\left(  \mathbf{\varphi,\theta},\mathbf{\varepsilon,}%
\eta\right)  =\mathcal{E}_{1C}\left(  \mathbf{\varphi}\right)  -\sum_{1\leq
j,k\leq2s}\varepsilon_{kj}\left(  \left\langle \left.  \varphi_{j}\right\vert
\varphi_{k}\right\rangle -n_{j}\mathbf{\left(  \mathbf{\theta}\right)  }%
\delta_{jk}\right)  -\eta\left\{  \sum_{1\leq j\leq2s}n_{j}\mathbf{\left(
\mathbf{\theta}\right)  }-1\right\}  \label{LI2}%
\end{equation}
respectively, where $\mathbf{n}$ is expressed in terms of the variable
$\mathbf{\theta}$ as
\begin{equation}
n_{j}\mathbf{\left(  \mathbf{\theta}\right)  }=\sin^{2}\theta_{j}.
\end{equation}
One carries out sequential optimizations over $\left(  \mathbf{\tilde{\varphi
},\tilde{\varepsilon}}\right)  $ and $\left(  \mathbf{\theta,}\eta\right)  ,$
giving
\begin{align}
\mathcal{L}\left(  \mathbf{\chi\left(  \mathbf{\theta}\right)  ,\epsilon
\left(  \mathbf{\theta}\right)  ,}\eta\right)   &  =\max_{\mathbf{\epsilon}%
}\min_{\mathbf{\chi}}\mathcal{L}\left(  \mathbf{\chi,\theta},\mathbf{\epsilon
,}\eta\right) \nonumber\\
\frac{\delta\mathcal{L}\left(  \mathbf{\chi\left(  \mathbf{\theta}\right)
,\mathbf{\theta},\epsilon\left(  \mathbf{\theta}\right)  ,}\eta\right)
}{\delta\mathbf{\chi}}  &  =\mathbf{0};\quad\frac{\partial\mathcal{L}\left(
\mathbf{\chi\left(  \mathbf{\theta}\right)  ,\mathbf{\theta},\epsilon\left(
\mathbf{\theta}\right)  ,}\eta\right)  }{\partial\mathbf{\epsilon}}=\mathbf{0}
\label{AGPL1}%
\end{align}
and
\begin{align}
\mathcal{L}\left(  \mathbf{\chi}\left(  \mathbf{\theta}_{\#}\left(  \eta
_{\#}\right)  \right)  \mathbf{,\epsilon}\left(  \mathbf{\theta}_{\#}\left(
\eta_{\#}\right)  \right)  \right)   &  =\max_{\mathbf{\eta}}\min
_{\mathbf{\theta}}\mathcal{L}\left(  \mathbf{\chi\left(  \mathbf{\theta
}\right)  ,\epsilon\left(  \mathbf{\theta}\right)  ,}\eta\right) \nonumber\\
\frac{\partial\mathcal{L}\left(  \mathbf{\chi}\left(  \mathbf{\theta}%
_{\#}\left(  \eta_{\#}\right)  \right)  ,\mathbf{\epsilon}\left(
\mathbf{\theta}_{\#}\left(  \eta_{\#}\right)  \right)  \mathbf{,}\eta
_{\#}\right)  }{\partial\mathbf{\theta}}  &  =\mathbf{0};\quad\frac
{\partial\mathcal{L}\left(  \mathbf{\chi}\left(  \mathbf{\theta}_{\#}\left(
\eta_{\#}\right)  \right)  ,\mathbf{\epsilon}\left(  \mathbf{\theta}%
_{\#}\left(  \eta_{\#}\right)  \right)  \mathbf{,}\eta_{\#}\right)  }%
{\partial\eta}=\mathbf{0} \label{AGPL2}%
\end{align}
for $\mathbf{\chi=\tilde{\varphi},\varphi}$ and $\mathbf{\epsilon
=\tilde{\varepsilon},\varepsilon.}$ The stationarity condition in eq. $\left(
\ref{AGPL1}\right)  $ gives
\begin{equation}
\frac{\delta\mathcal{E}_{1}\left(  \mathbf{\chi\left(  \mathbf{\theta}\right)
}\right)  }{\delta\chi_{j}}=\sum_{1\leq jk\leq2s}\epsilon\mathbf{\left(
\mathbf{\theta}\right)  }_{kj}\left\vert \varphi_{k}\mathbf{\left(
\mathbf{\theta}\right)  }\right\rangle ,
\end{equation}
which can be expressed \cite{WeinerOrtiz05} in terms of a generalized Fock
operator as
\begin{equation}
F\left(  \mathbf{\chi\left(  \mathbf{\theta}\right)  }\right)  \left\vert
\varphi_{j}\mathbf{\left(  \mathbf{\theta}\right)  }\right\rangle =\sum_{1\leq
jk\leq2s}\epsilon\mathbf{\left(  \mathbf{\theta}\right)  }_{kj}\left\vert
\varphi_{k}\mathbf{\left(  \mathbf{\theta}\right)  }\right\rangle ,
\end{equation}
while the stationarity condition eq. $\left(  \ref{AGPL2}\right)  $ gives
\begin{equation}
\frac{\partial\mathcal{L}\left(  \mathbf{\chi}\left(  \mathbf{\theta}%
_{\#}\left(  \eta_{\#}\right)  \right)  ,\mathbf{\epsilon}\left(
\mathbf{\theta}_{\#}\left(  \eta_{\#}\right)  \right)  \mathbf{,}\eta
_{\#}\right)  }{\partial\theta_{j}}=\frac{\partial\mathcal{L}\left(
\mathbf{\chi}\left(  \mathbf{\theta}_{\#}\left(  \eta_{\#}\right)  \right)
,\mathbf{\epsilon}\left(  \mathbf{\theta}_{\#}\left(  \eta_{\#}\right)
\right)  \mathbf{,}\eta_{\#}\right)  }{\partial n_{j}}\sin\left\{  2\theta
_{j}\right\}  .
\end{equation}
This leads to
\begin{equation}
\left\{  \frac{\partial\mathcal{E}_{1}\left(  \mathbf{\chi}\left(
\mathbf{n}\left(  \mathbf{\theta}_{\#}\left(  \eta_{\#}\right)  \right)
\right)  \right)  }{\partial n_{j}}-\eta_{\#}\right\}  \sin\left\{
2\theta_{j}\left(  \eta_{\#}\right)  \right\}  =0,
\end{equation}
which gives
\begin{equation}
\frac{\partial\mathcal{E}_{1}\left(  \mathbf{\chi}\left(  \mathbf{n}\left(
\mathbf{\theta}_{\#}\left(  \eta_{\#}\right)  \right)  \right)  \right)
}{\partial n_{j}}=\eta_{\#}%
\end{equation}
or
\begin{equation}
\sin\left\{  2\theta_{j}\left(  \eta_{\#}\right)  \right\}  =0\Rightarrow
\theta_{j}\left(  \eta_{\#}\right)  =0\text{ or }\frac{\pi}{2}\Rightarrow
n_{j}\left(  \mathbf{\theta}\left(  \eta_{\#}\right)  \right)  =0\text{ or }1.
\end{equation}
On the other hand sensitivity theory gives%
\begin{equation}
\frac{\partial\mathcal{E}_{1}\left(  \mathbf{\varphi}\left(  \mathbf{n}\left(
\mathbf{\theta}\right)  \right)  \right)  }{\partial n_{j}}=\varepsilon
_{jj}\mathbf{\left(  \mathbf{\theta}\right)  }%
\end{equation}
for all values of $\mathbf{\theta.}$

We are interested in the dependence on the occupation numbers of the
$2N$-electron state which we can obtain by using the chain rule
\begin{align}
&  \frac{\partial\mathcal{E}\left(  \mathbf{n,\varphi}\right)  }{\partial
N_{j}}=\sum_{1\leq k\leq2s}\frac{\partial\mathcal{E}\left(  \mathbf{n,\varphi
}\right)  }{\partial n_{k}}\frac{\partial n_{k}}{\partial N_{j}}\nonumber\\
&  \Rightarrow\frac{\partial\mathcal{E}\left(  \mathbf{n,\varphi}\right)
}{\partial N_{j}}=\sum_{1\leq k\leq2s}\frac{\partial\mathcal{E}\left(
\mathbf{n,\varphi}\right)  }{\partial n_{k}}\frac{\partial n_{k}}{\partial
N_{j}}=\sum_{1\leq k\leq2s}\frac{\partial\mathcal{E}\left(  \mathbf{n,\varphi
}\right)  }{\partial n_{k}}\left(  \frac{\partial N_{k}}{\partial n_{j}%
}\right)  ^{-1}%
\end{align}
i.e.%
\begin{equation}
\frac{\partial\mathcal{E}\left(  \mathbf{n,\varphi}\right)  }{\partial
\mathbf{N}}=\frac{\partial\mathcal{E}\left(  \mathbf{n,\varphi}\right)
}{\partial\mathbf{n}}\left(  \frac{\partial\mathbf{N}}{\partial\mathbf{n}%
}\right)  ^{-1}.
\end{equation}
The occupation numbers of the AGP\ state are explicitly given by
\begin{equation}
\mathbf{N}=\mathcal{N}\left(  \mathbf{n}\right)  ,
\end{equation}
where
\begin{equation}
N_{j}=\mathcal{N}_{j}\left(  \mathbf{n}\right)  =n_{j}\frac{S_{M-1}\left(
\jmath\right)  }{S_{M}}=n_{j}\frac{\partial\ln S_{M}}{\partial n_{j}}%
;\quad1\leq j\leq s.
\end{equation}
The matrix elements of the Jacobian $\frac{\partial\mathbf{N}}{\partial
\mathbf{n}}$ are
\begin{align}
\frac{\partial N_{k}}{\partial n_{j}}  &  =\frac{\partial}{\partial n_{j}%
}\left\{  n_{k}\frac{\partial\ln S_{M}}{\partial n_{k}}\right\}  =\delta
_{jk}\frac{\partial\ln S_{M}}{\partial n_{k}}+n_{k}\frac{\partial^{2}\ln
S_{M}}{\partial n_{j}\partial n_{k}}\nonumber\\
&  =\delta_{jk}\frac{S_{M-1}\left(  \jmath\right)  }{S_{M}}-n_{k}\frac
{S_{M-1}\left(  \jmath\right)  S_{M-1}\left(  k\right)  }{S_{M}^{2}}%
+n_{k}\frac{S_{M-2}\left(  jk\right)  }{S_{M}}.
\end{align}


\subsection{Unconstrained Variation of the AGP Energy Functional}

The occupation number variation nonunitary subgroup $G_{num}$ is given by
\begin{equation}
G_{num}=\left\{  \left.  k\left(  \mathbf{\eta}\right)  =\exp\left\{
\sum_{1\leq j\leq s}\eta_{j}\Omega_{j}\right\}  \right\vert \eta_{j}%
\in\mathbb{R}\right\}  , \label{occup}%
\end{equation}
where
\begin{equation}
\Omega_{j}=\left(  a_{j}^{\dagger}a_{j}+a_{j+s}^{\dagger}a_{j+s}\right)
;\quad1\leq j\leq s.
\end{equation}
The manifold of $2N$-electron AGP\ states $\mathcal{M}$ can be expressed as
\begin{equation}
\mathcal{M}\mathcal{=}\left\{  \left.  \overset{}{k\left(  \mathbf{\eta
}\right)  u\left(  \mathbf{\lambda}\right)  }\left\vert g_{0}^{N}\right\rangle
\right\vert \right.  \left.  \mathbf{\lambda}=\mathbf{\lambda}^{\dagger
}\right\}  .
\end{equation}
Explicitly the action on a reference geminal is given by
\begin{equation}
\left\vert \tilde{g}\left(  \mathbf{\eta},\mathbf{\lambda}\right)
\right\rangle =k\left(  \mathbf{\eta}\right)  u\left(  \mathbf{\lambda
}\right)  \left\vert g_{0}\right\rangle =\sum_{1\leq l\leq s}e^{\eta_{l}%
}c_{0l}\left\vert \varphi_{l}\left(  \mathbf{\lambda}\right)  \varphi
_{l+s}\left(  \mathbf{\lambda}\right)  \right\rangle ,
\end{equation}
where
\begin{equation}
\left\vert \varphi_{l}\left(  \mathbf{\lambda}\right)  \right\rangle
=\exp\left\{  i\sum_{1\leq j,k\leq2s}\lambda_{jk}a_{j}^{\dagger}a_{k}\right\}
\left\vert \varphi_{0l}\right\rangle .
\end{equation}
One can define a normalized geminal by
\begin{equation}
\left\vert g\left(  \mathbf{\eta},\mathbf{\lambda}\right)  \right\rangle
=\mathit{s}\left(  \mathbf{\eta}\right)  ^{-\frac{1}{2}}\left\vert \tilde
{g}\left(  \mathbf{\eta},\mathbf{\lambda}\right)  \right\rangle ,
\end{equation}
where
\begin{equation}
\mathit{s}\left(  \mathbf{\eta}\right)  =\left\langle \left.  \tilde{g}\left(
\mathbf{\eta},\mathbf{\lambda}\right)  \right\vert \tilde{g}\left(
\mathbf{\eta},\mathbf{\lambda}\right)  \right\rangle =\sum_{1\leq l\leq
s}e^{2\eta_{l}}c_{0l}^{2}=\sum_{1\leq l\leq s}e^{2\eta_{l}}n_{0l}%
\end{equation}
giving
\begin{align}
\left\vert g\left(  \mathbf{\eta},\mathbf{\lambda}\right)  \right\rangle  &
=\mathit{s}\left(  \mathbf{\eta}\right)  ^{-\frac{1}{2}}\sum_{1\leq l\leq
s}e^{\eta_{l}}c_{0l}\left\vert \varphi_{l}\left(  \mathbf{\lambda}\right)
\varphi_{l+s}\left(  \mathbf{\lambda}\right)  \right\rangle \nonumber\\
&  =\sum_{1\leq l\leq s}e^{\eta_{l}}c_{0l}\mathit{s}\left(  \mathbf{\eta
}\right)  ^{-\frac{1}{2}}\left\vert \varphi_{l}\left(  \mathbf{\lambda
}\right)  \varphi_{l+s}\left(  \mathbf{\lambda}\right)  \right\rangle
=\sum_{1\leq l\leq s}c_{l}\left(  \mathbf{\eta}\right)  \left\vert \varphi
_{l}\left(  \mathbf{\lambda}\right)  \varphi_{l+s}\left(  \mathbf{\lambda
}\right)  \right\rangle ,
\end{align}
where
\begin{equation}
c_{l}\left(  \mathbf{\eta}\right)  =e^{\eta_{l}}c_{0l}\mathit{s}\left(
\mathbf{\eta}\right)  ^{-\frac{1}{2}}%
\end{equation}
leading to
\begin{equation}
n_{l}\left(  \mathbf{\eta}\right)  =e^{2\eta_{l}}n_{0l}\mathit{s}\left(
\mathbf{\eta}\right)  ^{-1}.
\end{equation}
This gives
\begin{equation}
\sum_{1\leq l\leq s}n_{l}\left(  \mathbf{\eta}\right)  =\mathit{s}\left(
\mathbf{\eta}\right)  ^{-1}\sum_{1\leq l\leq s}e^{2\eta_{l}}n_{0l}=\left(
\sum_{1\leq l\leq s}e^{2\eta_{l}}n_{0l}\right)  ^{-1}\sum_{1\leq l\leq
s}e^{2\eta_{l}}n_{0l}=1
\end{equation}
which shows that $\left\{  n_{l}\left(  \mathbf{\eta}\right)  \right\}  $
satisfy the geminal occupation number positivity and normalization
constraints. We can define%
\begin{equation}
\left\vert \varphi_{l}\left(  \mathbf{\eta,\lambda}\right)  \right\rangle
=\exp\eta_{l}\exp\left\{  i\sum_{1\leq j,k\leq2s}\lambda_{jk}a_{j}^{\dagger
}a_{k}\right\}  \left\vert \varphi_{0l}\right\rangle ,
\end{equation}
so that
\begin{equation}
\left\langle \left.  \varphi_{l}\left(  \mathbf{\eta,\lambda}\right)
\right\vert \varphi_{l}\left(  \mathbf{\eta,\lambda}\right)  \right\rangle
=e^{2\eta_{l}}\left\langle \left.  \varphi_{0l}\right\vert \varphi
_{0l}\right\rangle =e^{2\eta_{l}}%
\end{equation}
and
\begin{equation}
\left\vert g\left(  \mathbf{\eta},\mathbf{\lambda}\right)  \right\rangle
=\mathit{s}\left(  \mathbf{\eta}\right)  ^{-\frac{1}{2}}\sum_{1\leq l\leq
s}c_{0l}\left\vert \varphi_{l}\left(  \mathbf{\eta,\lambda}\right)
\varphi_{l+s}\left(  \mathbf{\eta,\lambda}\right)  \right\rangle .
\end{equation}


We can express the geminal occupation numbers (and hence the occupation
numbers of the FORDM) and the normalization of the CGSO's in terms of
$\mathbf{\eta}$
\begin{align}
n_{l}\left(  \mathbf{\eta}\right)   &  =e^{2\eta_{l}}n_{0l}\mathit{s}\left(
\mathbf{\eta}\right)  ^{-1}\nonumber\\
s_{l}  &  =\left\langle \left.  \varphi_{l}\left(  \mathbf{\eta,\lambda
}\right)  \right\vert \varphi_{l}\left(  \mathbf{\eta,\lambda}\right)
\right\rangle =e^{2\eta_{l}}%
\end{align}
giving
\begin{equation}
\eta_{l}=\frac{1}{2}\ln\left\langle \left.  \varphi_{l}\left(  \mathbf{\eta
,\lambda}\right)  \right\vert \varphi_{l}\left(  \mathbf{\eta,\lambda}\right)
\right\rangle \equiv\frac{1}{2}\ln\left\langle \left.  \varphi_{l}\left(
\mathbf{\lambda}\right)  \right\vert \varphi_{l}\left(  \mathbf{\lambda
}\right)  \right\rangle =\frac{1}{2}\ln s_{l}%
\end{equation}
Hence%
\begin{equation}
\mathcal{E}\left(  \mathbf{n}_{\min}\mathbf{,\varphi}_{\min}\right)
=\mathcal{E}\left(  \mathbf{\eta}_{\min}\mathbf{,\lambda}_{\min}\right)
=\mathcal{E}\left(  \mathbf{\varphi}_{\min}\right)
\end{equation}
where
\begin{align}
\mathcal{E}\left(  \mathbf{n}_{\min}\mathbf{,\varphi}_{\min}\right)   &
=\min_{\mathbf{n\in\mathcal{S},\varphi\in}\mathcal{O}\left(  \mathbf{1}%
_{s}\right)  }\mathcal{E}\left(  \mathbf{n,\varphi}\right) \\
\mathcal{E}\left(  \mathbf{\eta}_{\min}\mathbf{,\lambda}_{\min}\right)   &
=\min_{\mathbf{\eta,\lambda}}\mathcal{E}\left(  \mathbf{\eta,\lambda}\right)
\\
\mathcal{E}\left(  \mathbf{\varphi}_{\min}\right)   &  =\min_{\mathbf{s\in
}\mathbb{R}_{+}^{s}}\min_{\mathbf{\varphi\in}\mathcal{O}\left(  \mathbf{s}%
\right)  }\mathcal{E}\left(  \mathbf{\varphi}\right)
\end{align}%
\begin{align}
\mathcal{S}  &  =\left\{  \left.  \mathbf{n}\right\vert n_{j}\geq0,1\leq j\leq
s;\sum_{1\leq j\leq s}n_{j}=1\right\} \\
\mathcal{O}\left(  \mathbf{s}\right)   &  =\left\{  \left.  \mathbf{\varphi
}\right\vert \left\langle \left.  \varphi_{j}\right\vert \varphi
_{k}\right\rangle =s_{j}\delta_{jk},1\leq j,k\leq s\right\}  .
\end{align}%
\begin{equation}
\mathbf{\lambda}=\mathbf{\lambda}^{\dagger},\mathbf{\eta\in}\mathbb{R}^{s}%
\end{equation}%
\begin{align}
\mathcal{L}\left(  \mathbf{\varphi,s}\right)   &  =\mathcal{E}\left(
\mathbf{n}_{0}\mathbf{,\varphi}\right)  -\sum_{1\leq j,k\leq2s}\varepsilon
_{kj}\left(  \left\langle \left.  \varphi_{j}\right\vert \varphi
_{k}\right\rangle -s_{j}\delta_{jk}\right) \\
\frac{\partial\mathcal{L}\left(  \mathbf{\varphi,s}\right)  }{\partial s_{j}}
&  =\varepsilon_{jj}.
\end{align}%
\begin{equation}
\mathcal{L}\left(  \mathbf{n\left(  \mathbf{\eta}\right)  ,\varphi}\right)
=\mathcal{L}\left(  \mathbf{\varphi,s}\left(  \mathbf{\eta}\right)  \right)
\end{equation}


\section{\bigskip Discussion}

\appendix

\section{Reduced Density Operators \label{Denmat}}

\bigskip

\section{Bases \label{Bases}}

\bigskip

\section{Operator Sets and Subspaces \label{Operators}}

\bigskip

\section{State Normalization \label{Normalization}}

If $D^{N}\in\mathcal{B}_{2}\left(  \mathcal{H}^{N}\right)  $ then
\begin{equation}
\operatorname{Tr}\left\{  \mathcal{N}D^{N}\right\}  =N\operatorname{Tr}%
\left\{  D^{N}\right\}  ,
\end{equation}
where
\begin{equation}
\mathcal{N}=\int\Phi\left(  \mathbf{x}\right)  ^{\dagger}\Phi\left(
\mathbf{x}\right)  d\mathbf{x,}%
\end{equation}
as
\begin{equation}
\mathcal{N}\left\vert \Psi\right\rangle =N\left\vert \Psi\right\rangle
\quad\forall\left\vert \Psi\right\rangle \in\mathcal{H}^{N}.
\end{equation}
For any state $D$
\begin{equation}
\operatorname{Tr}\left\{  \mathcal{N}D\right\}  =\int\gamma\left(
\mathbf{r,\xi}\right)  d\mathbf{r}d\mathbf{\xi,}%
\end{equation}
thus
\begin{equation}
N\operatorname{Tr}\left\{  D^{N}\right\}  =N\Rightarrow\operatorname{Tr}%
\left\{  D^{N}\right\}  =1.
\end{equation}


\section{Expansion in Terms of GSO's and Occupation Numbers\label{GSO&Occ}}

One can express the HK energy functional in terms of orthonormal GSO's
$\mathbf{\tilde{\varphi}}$ and occupation numbers $\mathbf{N}$ by the orbital
functional $\mathcal{\tilde{F}}_{O}$
\begin{equation}
\mathcal{\tilde{F}}_{O}\left(  \mathbf{\tilde{\varphi},N}\right)
=\mathcal{F}_{HK}\left(  \Gamma_{2}\left(  \mathbf{\tilde{\varphi},N}\right)
\right)  =\mathcal{F}_{HK}\left(  \gamma\right)  ,
\end{equation}
whose value is the same as the HK energy functional.

The functional $\mathcal{\tilde{F}}_{O}$ can be decomposed as
\begin{equation}
\mathcal{\tilde{F}}_{O}\left(  \mathbf{\tilde{\varphi},N}\right)
=\mathcal{\tilde{F}}_{T}\left(  \mathbf{\tilde{\varphi},N}\right)
+\mathcal{\tilde{F}}_{C}\left(  \mathbf{\tilde{\varphi},N}\right)
+\mathcal{\tilde{F}}_{OXC}\left(  \mathbf{\tilde{\varphi},N}\right)  ,
\end{equation}
where
\begin{equation}
\mathcal{\tilde{F}}_{T}\left(  \mathbf{\tilde{\varphi},N}\right)  =\sum_{1\leq
j\leq r_{1}}\left\langle \tilde{\varphi}_{j}\left\vert K\tilde{\varphi}%
_{j}\right.  \right\rangle N_{j},
\end{equation}
the orbital exchange-corellation functional is defined as
\begin{equation}
\mathcal{\tilde{F}}_{OXC}\left(  \mathbf{\tilde{\varphi},N}\right)
=\mathcal{F}_{XC}\left(  \Gamma_{2}\left(  \mathbf{\tilde{\varphi},N}\right)
\right)  +\mathcal{F}_{K}\left(  \Gamma_{2}\left(  \mathbf{\tilde{\varphi}%
,N}\right)  \right)  -\mathcal{\tilde{F}}_{T}\left(  \mathbf{\tilde{\varphi
},N}\right)  ,
\end{equation}
and one recalls the definitions
\begin{equation}
\mathcal{F}_{XC}\left(  \Gamma_{1}\left(  \mathbf{\varphi}\right)  \right)
=\operatorname{Tr}\left\{  V_{2}D_{\ast}\left(  \Gamma_{1}\left(
\mathbf{\varphi}\right)  \right)  \right\}  -\mathcal{F}_{C}\left(  \Gamma
_{1}\left(  \mathbf{\varphi}\right)  \right)
\end{equation}
and
\begin{equation}
\mathcal{F}_{K}\left(  \Gamma_{1}\left(  \mathbf{\varphi}\right)  \right)
=\operatorname{Tr}\left\{  KD_{\ast}\left(  \Gamma_{1}\left(  \mathbf{\varphi
}\right)  \right)  \right\}  .
\end{equation}


The minimization problem that determines the ground state energy and
space-spin density of a system characterized by a local potential $\mathrm{v}$
is
\begin{equation}
E\left(  \Gamma_{2}\left(  \mathbf{\tilde{\varphi}}_{GS}\mathbf{,N}\right)
\mathbf{,}\mathrm{v}\right)  =\min_{\left(  \mathbf{\tilde{\varphi},N}\right)
\mathbf{\in}\mathfrak{S}_{2}\left(  N\right)  \times\mathfrak{N}\left(
N\right)  }\mathcal{E}_{O}\left(  \mathbf{\tilde{\varphi},N,}\mathrm{v}%
\right)  , \label{KSminA}%
\end{equation}
where
\begin{equation}
\mathcal{E}_{O}\left(  \mathbf{\tilde{\varphi},N,}\mathrm{v}\right)
=\mathcal{\tilde{F}}_{O}\left(  \mathbf{\tilde{\varphi},N}\right)
+f_{\Gamma_{2}\left(  \mathbf{\tilde{\varphi},N}\right)  }\left(
\mathrm{v}\right)  .
\end{equation}
The feasible set $\mathfrak{S}_{2}\left(  N\right)  \times\mathfrak{N}\left(
N\right)  $ is generated by the constraints
\begin{equation}
\left\langle \left.  \tilde{\varphi}_{j}\right\vert \tilde{\varphi}%
_{k}\right\rangle =\delta_{jk},\,1\leq j,k\leq r_{1},
\end{equation}
and
\begin{equation}
0\leq N_{j}\leq1,\,1\leq j\leq r_{1};\sum_{\,1\leq j\leq r_{1}}N_{j}=N.
\end{equation}
The GSO minimization problem eq. $\left(  \ref{KSminA}\right)  $ can be
decomposed into two sequential optimizations, the first one determines the
saddle point $\mathcal{L}_{\mathbf{\tilde{\varphi}}}\left(  \mathbf{\tilde
{\varphi},N},\mathbf{\varepsilon},\mathrm{v}\right)  $ of the Lagrangian
\begin{equation}
\mathcal{L}_{\mathbf{\tilde{\varphi}}}\left(  \mathbf{\tilde{\varphi}%
},\mathbf{\tilde{\varepsilon}}\left(  \mathbf{N}\right)  \mathbf{,N,}%
\mathrm{v}\right)  =\mathcal{E}_{O}\left(  \mathbf{\tilde{\varphi}%
,N,}\mathrm{v}\right)  -\sum_{1\leq j,k\leq r_{1}}\tilde{\varepsilon}%
_{jk}\left\{  \overset{}{\left\langle \left.  \tilde{\varphi}_{j}\right\vert
\tilde{\varphi}_{k}\right\rangle -\delta_{jk}}\right\}  ,
\end{equation}
i.e.
\begin{equation}
\mathcal{L}_{\mathbf{\tilde{\varphi}}}\left(  \mathbf{\tilde{\varphi}}\left(
\mathbf{N}\right)  ,\mathbf{\tilde{\varepsilon}}\left(  \mathbf{N}\right)
,\mathbf{N,}\mathrm{v}\right)  =\max_{\mathbf{\tilde{\varepsilon}}}%
\min_{\mathbf{\tilde{\varphi}}}\mathcal{L}_{\mathbf{\tilde{\varphi}}}\left(
\mathbf{\tilde{\varphi}},\mathbf{\tilde{\varepsilon}}\left(  \mathbf{N}%
\right)  \mathbf{,N,}\mathrm{v}\right)  , \label{KSM1A}%
\end{equation}
which leads to optimal values $\mathcal{E}_{O}\left(  \mathbf{\mathbf{\tilde
{\varphi}}\left(  \mathbf{N}\right)  ,N,}\mathrm{v}\right)  =$ $\mathcal{L}%
_{\mathbf{\tilde{\varphi}}}\left(  \mathbf{\tilde{\varphi}}\left(
\mathbf{N}\right)  ,\mathbf{\tilde{\varepsilon}}\left(  \mathbf{N}\right)
,\mathbf{N,}\mathrm{v}\right)  $ of the GSO energy functional for fixed values
of the control parameters $\mathbf{N}$. In the second optimization one finds
minima of $\mathcal{E}_{O}\left(  \mathbf{\mathbf{\tilde{\varphi}}\left(
\mathbf{N}\right)  ,N,}\mathrm{v}\right)  $ with respect the variables
$\mathbf{N}$ that satisfy the constraints
\begin{equation}
0\leq N_{j}\leq1,1\leq j\leq r_{1}%
\end{equation}
which can be held by setting
\begin{equation}
N_{j}\left(  \mathbf{\theta}\right)  =\sin^{2}\theta_{j},
\end{equation}
and the introduction of a Lagrange parameter $\eta$ in order to enforce the
constraint
\begin{equation}
\sum_{1\leq j\leq r_{1}}N_{j}\left(  \mathbf{\theta}\right)  =N.
\end{equation}
This leads to a Lagrangian
\begin{equation}
\mathcal{L}_{\mathbf{\theta}}\left(  \mathbf{\tilde{\varphi}}\left(
\mathbf{\theta}\right)  ,\mathbf{\theta},N,\eta,\mathrm{v}\right)
=\mathcal{E}_{O}\left(  \mathbf{\tilde{\varphi}}\left(  \mathbf{\theta
}\right)  \mathbf{,\theta},\mathrm{v}\right)  -\eta\left\{  \sum_{\,1\leq
j\leq r_{1}}N_{j}\left(  \mathbf{\theta}\right)  -N\right\}  .
\end{equation}
The determination of the saddle point $\mathcal{L}_{\mathbf{\theta}}\left(
\mathbf{\tilde{\varphi}}\left(  \mathbf{\theta}_{\#}\left(  N\right)  \right)
,\mathbf{\theta}_{\#}\left(  N\right)  ,\eta\left(  N\right)  ,\mathrm{v}%
\right)  $ given by
\begin{equation}
\mathcal{L}_{\mathbf{\theta}}\left(  \mathbf{\tilde{\varphi}}\left(
\mathbf{\theta}_{\#}\left(  N\right)  \right)  ,\mathbf{\theta}_{\#}\left(
N\right)  ,\eta\left(  N\right)  ,\mathrm{v}\right)  =\max_{\eta}%
\min_{\mathbf{\theta}}\mathcal{L}_{\mathbf{\theta}}\left(  \mathbf{\tilde
{\varphi}}\left(  \mathbf{\theta}\right)  ,\mathbf{\theta},N,\eta
,\mathrm{v}\right)  .
\end{equation}
determines the ground state energy $E\left(  \gamma_{GS}\mathbf{,}%
\mathrm{v}\right)  =\mathcal{E}_{O}\left(  \mathbf{\tilde{\varphi}}\left(
\mathbf{\theta}_{\#}\left(  N\right)  \right)  ,\mathbf{\theta}_{\#}\left(
N\right)  ,\eta,N,\mathrm{v}\right)  .$The ground state density can then
expanded as
\begin{equation}
\gamma_{GS}\left(  \mathbf{x}\right)  =\sum_{1\leq j\leq r_{1}}N_{j}\left\vert
\left[  \tilde{\varphi}_{j}\left(  \mathbf{\theta}_{\#}\left(  N\right)
\right)  \right]  \left(  \mathbf{x}\right)  \right\vert ^{2}.
\end{equation}
The saddle points of $\mathcal{L}_{\mathbf{\tilde{\varphi}}}\left(
\mathbf{\tilde{\varphi}},\mathbf{\tilde{\varepsilon},N,}\mathrm{v}\right)  $
are characterized by the necessary conditions
\begin{subequations}
\begin{align}
\frac{\delta\mathcal{L}_{\mathbf{\tilde{\varphi}}}\left(  \mathbf{\tilde
{\varphi}},\mathbf{\tilde{\varepsilon},N,}\mathrm{v}\right)  }{\delta
\varphi_{j}^{\ast}}  &  =0;\,1\leq j\leq r_{1},\label{KSC1A}\\
\frac{\partial\mathcal{L}_{\mathbf{\tilde{\varphi}}}\left(  \mathbf{\tilde
{\varphi}},\mathbf{\tilde{\varepsilon},N,}\mathrm{v}\right)  }{\partial
\tilde{\varepsilon}_{jk}}  &  =0;\,1\leq j,k\leq r_{1} \label{KSC2A}%
\end{align}
(Appendix \ref{ComplexConjugate} describes why it is sufficient to consider
derivatives wrt $\varphi_{j}^{\ast}).$ The condition eq. $\left(
\ref{KSC1A}\right)  $ implies that at an optimal point the matrix of Lagrange
multipliers $\mathbf{\tilde{\varepsilon}}\left(  \mathbf{N}\right)  $ is
hermitian i.e.
\end{subequations}
\begin{equation}
\tilde{\varepsilon}\left(  \mathbf{N}\right)  _{jk}=\tilde{\varepsilon}\left(
\mathbf{N}\right)  _{kj}^{\ast},\quad1\leq j,k\leq r_{1}%
\end{equation}
and
\begin{equation}
\frac{\delta\mathcal{E}_{O}\left(  \mathbf{\tilde{\varphi}}\left(
\mathbf{N}\right)  ,\mathbf{N,}\mathrm{v}\right)  }{\delta\tilde{\varphi}%
_{j}^{\ast}}-\sum_{1\leq k\leq r_{1}}\tilde{\varepsilon}_{jk}\left(
\mathbf{N}\right)  \left\vert \tilde{\varphi}_{k}\left(  \mathbf{N}\right)
\right\rangle =0;\,1\leq j\leq r_{1}.
\end{equation}
Which after noting that
\begin{align}
\frac{\delta\mathcal{E}_{O}\left(  \mathbf{\tilde{\varphi}}\left(
\mathbf{N}\right)  ,\mathbf{N,}\mathrm{v}\right)  }{\delta\tilde{\varphi}%
_{j}^{\ast}}  &  =\int\frac{\delta\mathcal{E}_{O}\left(  \mathbf{\tilde
{\varphi}}\left(  \mathbf{N}\right)  ,\mathbf{N,}\mathrm{v}\right)  }%
{\delta\gamma\left(  \mathbf{x}^{\prime}\right)  }\frac{\delta\gamma\left(
\mathbf{x}^{\prime}\right)  }{\delta\tilde{\varphi}_{j}^{\ast}}d\mathbf{x}%
^{\prime}\nonumber\\
&  =\int\frac{\delta\mathcal{E}_{O}\left(  \mathbf{\tilde{\varphi}}\left(
\mathbf{N}\right)  ,\mathbf{N,}\mathrm{v}\right)  }{\delta\gamma\left(
\mathbf{x}\right)  }\tilde{\varphi}_{j}\left(  \mathbf{x}^{\prime}\right)
N_{j}\delta\left(  \mathbf{x-x}^{\prime}\right)  d\mathbf{x}^{\prime}%
=\frac{\delta\mathcal{E}_{O}\left(  \mathbf{\tilde{\varphi}}\left(
\mathbf{N}\right)  ,\mathbf{N,}\mathrm{v}\right)  }{\delta\gamma\left(
\mathbf{x}\right)  }N_{j}\tilde{\varphi}_{j}\left(  \mathbf{x}\right)
\end{align}
gives%
\begin{equation}
\frac{\delta\mathcal{E}_{O}\left(  \mathbf{\tilde{\varphi}}\left(
\mathbf{N}\right)  ,\mathbf{N,}\mathrm{v}\right)  }{\delta\gamma}%
N_{j}\left\vert \tilde{\varphi}_{j}\left(  \mathbf{N}\right)  \right\rangle
-\sum_{1\leq k\leq r_{1}}\tilde{\varepsilon}_{jk}\left(  \mathbf{N}\right)
\left\vert \tilde{\varphi}_{k}\left(  \mathbf{N}\right)  \right\rangle
=0;\,1\leq j\leq r_{1}.
\end{equation}
The operator $\frac{\delta\mathcal{E}_{O}\left(  \mathbf{\tilde{\varphi}%
}\left(  \mathbf{N}\right)  ,\mathbf{N,}\mathrm{v}\right)  }{\delta\gamma}$ is
a multiplication operator given by
\begin{equation}
\frac{\delta\mathcal{E}_{O}\left(  \mathbf{\tilde{\varphi}}\left(
\mathbf{N}\right)  ,\mathbf{N,}\mathrm{v}\right)  }{\delta\gamma}\left\vert
\tilde{\varphi}_{j}\left(  \mathbf{N}\right)  \right\rangle =\left\vert
\frac{\delta\mathcal{E}_{O}\left(  \mathbf{\tilde{\varphi}}\left(
\mathbf{N}\right)  ,\mathbf{N,}\mathrm{v}\right)  }{\delta\gamma}%
\tilde{\varphi}_{j}\left(  \mathbf{N}\right)  \right\rangle ,
\end{equation}
where
\begin{equation}
\left[  \frac{\delta\mathcal{E}_{O}\left(  \mathbf{\tilde{\varphi}}\left(
\mathbf{N}\right)  ,\mathbf{N,}\mathrm{v}\right)  }{\delta\gamma}%
\tilde{\varphi}_{j}\left(  \mathbf{N}\right)  \right]  \left(  \mathbf{x}%
\right)  =\frac{\delta\mathcal{E}_{O}\left(  \mathbf{\tilde{\varphi}}\left(
\mathbf{N}\right)  ,\mathbf{N,}\mathrm{v}\right)  }{\delta\gamma\left(
\mathbf{x}\right)  }\left[  \tilde{\varphi}_{j}\left(  \mathbf{N}\right)
\right]  \left(  \mathbf{x}\right)  .
\end{equation}
We thus obtain the intermediate GSO equations that correspond to condition eq.
$\left(  \ref{KSC1A}\right)  $
\begin{equation}
\frac{\delta\mathcal{E}_{O}\left(  \mathbf{\tilde{\varphi}}\left(
\mathbf{N}\right)  ,\mathbf{N,}\mathrm{v}\right)  }{\delta\gamma}\left\vert
\tilde{\varphi}_{j}\left(  \mathbf{N}\right)  \right\rangle =\sum_{1\leq k\leq
r_{1}}\tilde{\varepsilon}_{jk}\left(  \mathbf{N}\right)  \left\vert
\tilde{\varphi}_{k}\left(  \mathbf{N}\right)  \right\rangle .
\end{equation}
By partial differentiation the rate of change of the orbital functional energy
with respect to changes in occupation numbers is
\begin{equation}
\frac{\partial\mathcal{E}_{O}\left(  \mathbf{\tilde{\varphi}}\left(
\mathbf{N}\right)  ,\mathbf{N,}\mathrm{v}\right)  }{\partial N_{j}}%
=\tilde{\varepsilon}_{jj}\left(  \mathbf{N}\right)  ,
\end{equation}
which must be equal to the result obtained through sensitivity theory when the
occupation numbers are treated as control parameters.

\section{Complex Conjugate Variables \label{ComplexConjugate}}

\section{Symmetric Function Relationships}%

\begin{equation}
\frac{\partial S_{N}}{\partial n_{j}}=S_{N-1}\left(  j\right)
\end{equation}%
\begin{align}
\sum_{j}n_{j}S_{N-1}\left(  j\right)   &  =NS_{N}\\
\sum_{k}n_{k}S_{N-2}\left(  jk\right)   &  =\left(  N-1\right)  S_{N-1}\left(
j\right) \\
\sum_{j<k}n_{j}n_{k}S_{N-2}\left(  jk\right)   &  =\frac{N\left(  N-1\right)
}{2}S_{N}%
\end{align}


For a $2N$-electron IPS the geminal has the form
\begin{equation}
\left\vert g\right\rangle =N^{-\frac{1}{2}}\sum_{1\leq j\leq N}\left\vert
\psi_{j}\psi_{j+s}\right\rangle
\end{equation}
which for a $8$-electron case is
\begin{equation}
\left\vert g\right\rangle =4^{-\frac{1}{2}}\sum_{1\leq j\leq4}\left\vert
\psi_{j}\psi_{j+s}\right\rangle =\frac{1}{2}\sum_{1\leq j\leq4}\left\vert
\psi_{j}\psi_{j+s}\right\rangle
\end{equation}
so that
\begin{equation}
n_{j}=\left\{
\begin{array}
[c]{c}%
N^{-1};1\leq j\leq N\\
\qquad0;\text{ }N+1\leq j\leq s
\end{array}
\right.
\end{equation}
which for a $6$-electron case is%
\begin{equation}
n_{j}=\left\{
\begin{array}
[c]{c}%
\frac{1}{4};1\leq j\leq4\\
0;\text{ }5\leq j\leq s
\end{array}
\right.
\end{equation}%
\begin{equation}
S_{N}=N^{-N}%
\end{equation}
which for a $8$-electron case is%
\begin{equation}
S_{4}=\frac{1}{256}.
\end{equation}
The first derivatives of the symmetric function are
\begin{equation}
S_{N-1}\left(  \jmath\right)  =\sum_{\substack{j\notin\left\{  j_{1}%
,\cdots,j_{N-1}\right\}  \\1\leq j_{1}<\cdots<j_{N-1}\leq N}}n_{j_{1}}\cdots
n_{j_{N-1}}=\left\{
\begin{array}
[c]{c}%
\binom{N}{N-1}N^{-\left(  N-1\right)  };\quad j\notin\left\{  1,\cdots
,N\right\} \\
\binom{N-1}{N-1}N^{-\left(  N-1\right)  };\quad j\in\left\{  1,\cdots
,N\right\}
\end{array}
\right.
\end{equation}
which for a $8$-electron $s=8$ case where $\left(  n_{1},n_{2},n_{3}%
,n_{4}\right)  =\left(  \frac{1}{4},\frac{1}{4},\frac{1}{4},\frac{1}%
{4}\right)  $ is%

\begin{align}
S_{3}\left(  8\right)   &  =n_{1}n_{2}n_{3}+n_{1}n_{2}n_{4}+n_{1}n_{2}%
n_{5}+n_{1}n_{2}n_{6}+n_{1}n_{2}n_{7}\\
&  +n_{1}n_{3}n_{4}+n_{1}n_{3}n_{5}+n_{1}n_{3}n_{6}+n_{1}n_{3}n_{7}\\
&  +n_{1}n_{4}n_{5}+n_{1}n_{4}n_{6}+n_{1}n_{4}n_{7}\\
&  +n_{1}n_{5}n_{6}+n_{1}n_{5}n_{7}+n_{1}n_{6}n_{7}\\
&  +n_{2}n_{3}n_{4}+n_{2}n_{3}n_{5}+n_{2}n_{3}n_{6}+n_{2}n_{3}n_{7}\\
&  +n_{2}n_{4}n_{5}+n_{2}n_{4}n_{6}+n_{2}n_{4}n_{7}\\
&  +n_{2}n_{5}n_{6}+n_{2}n_{5}n_{7}+n_{2}n_{6}n_{7}\\
&  +n_{3}n_{4}n_{5}+n_{3}n_{4}n_{6}+n_{3}n_{4}n_{7}\\
&  +n_{3}n_{5}n_{6}+n_{3}n_{5}n_{7}+n_{3}n_{6}n_{7}\\
&  +n_{4}n_{5}n_{6}+n_{4}n_{5}n_{7}+n_{4}n_{6}n_{7}+n_{5}n_{6}n_{7}%
\end{align}%
\begin{equation}
S_{4}\left(  \jmath\right)  =\sum_{\substack{j\notin\left\{  j_{1},j_{2}%
,j_{3},j_{4}\right\}  \\1\leq j_{1}<j_{2}<j_{3}<j_{4}\leq8}}n_{j_{1}}n_{j_{2}%
}n_{j_{3}}n_{j_{4}}=\left\{
\begin{array}
[c]{c}%
\frac{1}{16};\quad j\notin\left\{  1,2,3,4\right\} \\
\frac{1}{64};\quad j\in\left\{  1,2,3,4\right\}
\end{array}
\right.
\end{equation}%
\begin{equation}
S_{N-2}\left(  jk\right)  =\sum_{\substack{j,k\notin\left\{  j_{1}%
,\cdots,j_{N-2}\right\}  \\1\leq j_{1}<\cdots<j_{N-2}\leq N}}n_{j_{1}}\cdots
n_{j_{N-2}}=\left\{
\begin{array}
[c]{c}%
\binom{N}{N-2}N^{-\left(  N-2\right)  };\quad j,k\notin\left\{  1,\cdots
,N\right\} \\
\binom{N-1}{N-2}N^{-\left(  N-2\right)  };\quad j\notin\left\{  1,\cdots
,N\right\}  ,k\in\left\{  1,\cdots,N\right\} \\
\binom{N-1}{N-2}N^{-\left(  N-2\right)  };\quad k\notin\left\{  1,\cdots
,N\right\}  ,j\in\left\{  1,\cdots,N\right\} \\
\binom{N-2}{N-2}N^{-\left(  N-2\right)  };\quad j\in\left\{  1,\cdots
,N\right\}  ,k\in\left\{  1,\cdots,N\right\}
\end{array}
\right.
\end{equation}
which for a $8$-electron case is
\begin{align}
S_{2}\left(  7,8\right)  =n_{1}n_{2}  &  +n_{1}n_{3}+n_{1}n_{4}+n_{1}%
n_{5}+n_{1}n_{6}+n_{2}n_{3}+n_{2}n_{4}+n_{2}n_{5}\\
&  +n_{2}n_{6}+n_{3}n_{4}+n_{3}n_{5}+n_{3}n_{6}+n_{4}n_{5}+n_{4}n_{6}%
+n_{5}n_{6}%
\end{align}%
\begin{equation}
S_{2}\left(  jk\right)  =\sum_{\substack{j,k\notin\left\{  j_{1}%
,\cdots,j_{N-2}\right\}  \\1\leq j_{1}<\cdots<j_{N-2}\leq N}}n_{j_{1}}\cdots
n_{j_{N-2}}=\left\{
\begin{array}
[c]{c}%
\frac{3}{8};\quad j,k\notin\left\{  1,\cdots,N\right\}  \qquad\qquad\qquad\\
\frac{3}{16};\quad j\notin\left\{  1,\cdots,N\right\}  ,k\in\left\{
1,\cdots,N\right\} \\
\frac{3}{16};\quad k\notin\left\{  1,\cdots,N\right\}  ,j\in\left\{
1,\cdots,N\right\} \\
\frac{1}{16};\quad j\in\left\{  1,\cdots,N\right\}  ,k\in\left\{
1,\cdots,N\right\}
\end{array}
\right.
\end{equation}
The diagonal
\begin{equation}
\frac{\partial N_{k}}{\partial n_{k}}=\left\{
\begin{array}
[c]{c}%
N;\quad k\notin\left\{  1,\cdots,N\right\} \\
0;\quad k\in\left\{  1,\cdots,N\right\}
\end{array}
\right.
\end{equation}%
\begin{equation}
\frac{\partial N_{k}}{\partial n_{j}}=\left\{
\begin{array}
[c]{c}%
-n_{k}\frac{N\left(  N+1\right)  }{2};\quad j\neq k\notin\left\{
1,\cdots,N\right\} \\
n_{k};\quad j\notin\left\{  1,\cdots,N\right\}  ,k\in\left\{  1,\cdots
,N\right\} \\
n_{k};\quad k\notin\left\{  1,\cdots,N\right\}  ,j\in\left\{  1,\cdots
,N\right\} \\
0;\quad j\in\left\{  1,\cdots,N\right\}  ,k\in\left\{  1,\cdots,N\right\}
\end{array}
\right.
\end{equation}



\end{document}